% !TeX program = xelatex
\documentclass[12pt,a4paper]{article}

% =========================================================
% ЯЗЫК И ШРИФТЫ
% =========================================================

\usepackage{fontspec}
\usepackage{polyglossia}
\usepackage{dsfont}

\setdefaultlanguage{russian}
\setotherlanguage{english}

\setmainfont{CMU Serif}
\setsansfont{CMU Sans Serif}
\setmonofont{CMU Typewriter Text}


% =========================================================
% МАТЕМАТИКА
% =========================================================
\usepackage{amsmath}
\usepackage{amssymb}
\usepackage{amsfonts}
\usepackage{amsthm}
\usepackage{mathtools}
\usepackage{bm}
\usepackage{mathrsfs}

% =========================================================
% СТРАНИЦА
% =========================================================
\usepackage[a4paper, margin=2cm]{geometry}

% =========================================================
% ЦВЕТА И РАМКИ
% =========================================================
\usepackage{xcolor}
\usepackage[most]{tcolorbox}

% =========================================================
% ГРАФИКА
% =========================================================
\usepackage{graphicx}

% =========================================================
% СПИСКИ
% =========================================================
\usepackage{enumitem}

% =========================================================
% ЗАГОЛОВКИ
% =========================================================
\usepackage{titlesec}

% =========================================================
% КОЛОНТИТУЛЫ
% =========================================================
\usepackage{fancyhdr}

% =========================================================
% ТАБЛИЦЫ
% =========================================================
\usepackage{booktabs}
\usepackage{array}

% =========================================================
% ССЫЛКИ
% =========================================================
\usepackage{hyperref}

\hypersetup{
    colorlinks=true,
    linkcolor=blue,
    urlcolor=blue,
    citecolor=blue
}

% =========================================================
% НАСТРОЙКИ
% =========================================================
\setlength{\parindent}{0pt}
\setlength{\parskip}{6pt}

% -------------------------
% Определение
% -------------------------
\newtcolorbox{definition}[1][]{
    breakable,
    colback=red!5,
    colframe=red!75!black,
    fonttitle=\bfseries,
    title=Определение,
    #1
}

% -------------------------
% Теорема
% -------------------------
\newtcolorbox{theorem}[2][]{
    breakable,
    colback=green!5,
    colframe=green!50!black,
    fonttitle=\bfseries,
    title=Теорема #2,
    #1
}

% -------------------------
% Лемма
% -------------------------
\newtcolorbox{lemma}[2][]{
    breakable,
    colback=orange!5,
    colframe=orange!70!black,
    fonttitle=\bfseries,
    title=Лемма #2,
    #1
}

% -------------------------
% Следствие
% -------------------------
\newtcolorbox{corollary}[1][]{
    breakable,
    colback=green!5,
    colframe=green!35!black,
    fonttitle=\bfseries,
    title=Следствие,
    #1
}

% -------------------------
% Доказательство
% -------------------------
\newtcolorbox{proofbox}[2][]{
    breakable,
    colback=blue!5,
    colframe=blue!50,
    fonttitle=\bfseries,
    title=Доказательство #2,
    #1
}

% -------------------------
% Замечание
% -------------------------
\newtcolorbox{note}[1][]{
    breakable,
    colback=purple!5,
    colframe=purple!80!black,
    fonttitle=\bfseries,
    title=Замечание,
    #1
}

% -------------------------
% Свойства
% -------------------------
\newtcolorbox{props}[2][]{
    breakable,
    colback=green!5,
    colframe=green!70!black,
    fonttitle=\bfseries,
    title=Свойства #2,
    #1
}

% -------------------------
% Примеры
% -------------------------
\newtcolorbox{example}[1][]{
    breakable,
    colback=brown!5,
    colframe=brown!40!black,
    fonttitle=\bfseries,
    title=Примеры,
    #1
}

% =========================================================
% ДОКУМЕНТ
% =========================================================
\begin{document}

\tableofcontents

\newpage

\section{Алгебра\, $\sigma$-алгебра и мера}

\subsection{Алгебра и $\sigma$-алгебра}

Пусть X — унивёрсум, $A \subset X$

\begin{definition}

\textbf{Алгебра} — непустая система множеств ($\mathcal{M}$) из X, удовлетворяющая двум условиям:

1) Если множетсво лежит в $\mathcal{M}$, то его дополнение тоже лежит в $\mathcal{M}$

$A \in \mathcal{M} \Longrightarrow A^C \in \mathcal{M}$

2) Если n множеств лежит в $\mathcal{M}$, то их объединение тоже лежит в $\mathcal{M}$

$\forall n < \infty: A_1, A_2,...,A_n \in \mathcal{M} \Longrightarrow \cup_{i=1}^{n} A_i \in \mathcal{M}$

\end{definition}

\begin{note}
Если второе условие верно для счётного числа множеств, то $\mathcal{M}$ — $\sigma$-алгебра
\end{note}

\begin{example}

1) $\{\emptyset, X\}$

2) $2^X$

3) $X = \{1, 2, 3\};\quad \mathcal{M} = \{\emptyset, {1}, {2, 3}, X\}$

\end{example}

\begin{props}{алгебры}

1) $\emptyset, X\ \in \mathcal{M}$


2) $A, B \in \mathcal{M} \Longrightarrow A \cap B \in \mathcal{M}$


3) $A, B \in \mathcal{M} \Longrightarrow A \backslash B \in \mathcal{M}$


4) $A, B \in \mathcal{M} \Longrightarrow A \triangle B \in \mathcal{M}$


5) Если $\mathcal{M}$ — $\sigma$-алгебра, то объединение счётного количества множеств, лежащх в $\mathcal{M}$, лежит в $\mathcal{M}$
(Если $A_i \in \mathcal{M}$, то $\cup_{i=1}^{\infty} A_i \in \mathcal{M}$)

\end{props}

\begin{proofbox}{свойства 1}
Пусть $A \in \mathcal{M}$. Тогда $A^C \in \mathcal{M}$. Тогда $(A \cup A^C) = X \in \mathcal{M}$
\end{proofbox}

\begin{proofbox}{свойства 2}

\[
A, B \in \mathcal{M} \Longrightarrow A^C, B^C \in \mathcal{M} \Longrightarrow (A^C \cup B^C) \in \mathcal{M} \Longrightarrow (A^C \cup B^C)^C = (A \cap B) \in \mathcal{M}
\]

\end{proofbox}

\begin{proofbox}{свойства 3}

$B \in \mathcal{M} \Longrightarrow B^C \in \mathcal{M} \\$
Тогда $A \backslash B = A \cap B^C \in \mathcal{M}$

\end{proofbox}



\begin{proofbox}{свойства 4}

$A, B \in \mathcal{M} \Longrightarrow A \backslash B \in \mathcal{M}$ и $B \backslash A \in \mathcal{M}$

Тогда $A \backslash B \cup B \backslash A = A \triangle B \in \mathcal{M}$

\end{proofbox}


\subsection{Мера}

Пусть $A \cap B = \emptyset$. Тогда $A \sqcup B$ — \textbf{дизъюнктное объединение}

Пусть $\mathcal{M}$ — система множеств, $\emptyset \in \mathcal{M}$. Тогда

\begin{definition}

\textbf{Мера} — отображение $\mu: \mathcal{M} \rightarrow [0, +\infty]$, если:

1) $\mu(\emptyset) = 0$

2) $A_i \in \mathcal{M}, \sqcup_{i=1}^\infty A_i \in \mathcal{M} \Longrightarrow \mu(\sqcup_{i=1}^\infty A_i) = \sum_{i=1}^\infty \mu(A_i) \quad (\sigma-$аддитивность)

\end{definition}

\begin{example}

1) Нулевая мера: $\mu(A) = 0 \quad \forall A$
\begin{flalign*}
2)\quad \mu(A) &=
\begin{cases}
|A|, & \text{если A конечно},\\
\infty, & \text{иначе}
\end{cases}
&
\end{flalign*}

3) Дискретная мера: пусть $a \in X$. Тогда:
\begin{flalign*}
\mu(A) &=
\begin{cases}
1, & \text{если $a \in A$},\\
0, & \text{иначе}
\end{cases}
&
\end{flalign*}

\end{example}

Пусть $\mathcal{M}$ --- $\sigma$-алгебра, $\mu$ --- мера на $\mathcal{M}$, $A,B\in\mathcal{M}$.

\begin{props}{меры}

1. \textbf{Монотонность:}

$$
A\subset B \Longrightarrow \mu(A)\leq\mu(B)
$$


2. \textbf{Непрерывность.} Пусть $\mathcal{M}$ --- $\sigma$-алгебра, $\mu$ --- мера. Тогда:

Пусть $A_i\in\mathcal{M}$,
\(    A_1\subset A_2\subset\ldots\subset A_n\subset A_{n+1}\subset\ldots
   \)
Тогда
\(    \lim_{n\to\infty}\mu(A_n)
    =
    \mu\left(\bigsqcup_{i=1}^{\infty}A_i\right).
   \)


Пусть $A_i\in\mathcal{M}$,
\[
A_1\supset A_2\supset A_3\supset\ldots\supset A_n\supset A_{n+1}\supset\ldots
\]
и $\mu(A_1)<+\infty$. Тогда
\[
\lim_{n\to\infty}\mu(A_n)
=
\mu\left(\bigcap_{i=1}^{\infty}A_i\right)
\]

\end{props}

\begin{proofbox}{свойства 1}

Пусть $C=B\setminus A$. Тогда

$$
\mu(A\sqcup C)=\mu(A)+\mu(C)=\mu(B)
$$

откуда

$$
\mu(B)\geq\mu(A)
$$

\end{proofbox}

\begin{proofbox}{свойства 2}
Пусть $A_0=\varnothing$. Тогда
\[
\bigsqcup_{i=1}^{\infty}A_i
=
A_1\sqcup(A_2\setminus A_1)\sqcup(A_3\setminus A_2)\sqcup\ldots
\]
и, следовательно,
\[
\mu\left(\bigsqcup_{i=1}^{\infty}A_i\right)
=
\sum_{i=1}^{\infty}\mu(A_i\setminus A_{i-1})
\]
Поэтому
\[
\mu\left(\bigsqcup_{i=1}^{\infty}A_i\right)
=
\lim_{n\to\infty}
\sum_{i=1}^{n}\mu(A_i\setminus A_{i-1})
=
\lim_{n\to\infty}\mu(A_n)
\]
\end{proofbox}

\begin{proofbox}{свойства 3}
Обозначим
\[
B_n=A_n\setminus A_{n+1}
\]
Тогда
\[
A_1=
\left(\bigcap_{i=1}^{\infty}A_i\right)
\sqcup
\bigsqcup_{i=1}^{\infty}B_i
\]
Поэтому
\[
\mu(A_1)
=
\mu\left(\bigcap_{i=1}^{\infty}A_i\right)
+
\sum_{i=1}^{\infty}\mu(B_i)
\]
Для каждого $n$ имеем
\[
A_n=
\left(\bigcap_{i=1}^{\infty}A_i\right)
\sqcup
\bigsqcup_{i=n}^{\infty}B_i
\]
откуда
\[
\mu(A_n)
=
\mu\left(\bigcap_{i=1}^{\infty}A_i\right)
+
\sum_{i=n}^{\infty}\mu(B_i)
\]
Так как $\mu(A_1)<+\infty$, хвост ряда стремится к нулю. Следовательно,
\[
\lim_{n\to\infty}\mu(A_n)
=
\mu\left(\bigcap_{i=1}^{\infty}A_i\right)
\]
\end{proofbox}


\newpage

\section{Мера Лебега в $\mathbb{R}^n$}

\subsection{Определение}

Пусть $a,b\in\mathbb{R}^n$:
\[
a=(a^1,\ldots,a^n),\qquad b=(b^1,\ldots,b^n),\qquad a^i<b^i
\]

\[
(a,b)=(a^1,b^1)\times\ldots\times(a^n,b^n) \text{ — открытый брус}
\]

\[
[a,b]=[a^1,b^1]\times\ldots\times[a^n,b^n] \text{ — закрытый брус}
\]

\begin{definition}
\textbf{Брусом в $\mathbb{R}^n$} называется множество вида
\[
\Pi=\langle a,b\rangle
\]
где $a,b\in\mathbb{R}^n$, $a^i<b^i$.

\[
\mu(\Pi)=(b^1-a^1)(b^2-a^2)\cdots(b^n-a^n)
=\prod_{i=1}^n(b^i-a^i)
\]

Это называется \textbf{мерой Лебега} на брусах.
\end{definition}

\begin{lemma}{}
Пусть $P_k$ и $P$ --- брусы.

\[
P_k\in\mathcal P,\qquad P\in\mathcal P
\]

Тогда:

\[
\bigcup_{k=1}^{m}P_k\subset P \Longrightarrow \sum_{k=1}^{m}\mu(P_k)\leq\mu(P)
\]

\[
P\subset\bigcup_{k=1}^{m}P_k \Longrightarrow \mu(P)\leq\sum_{k=1}^{m}\mu(P_k)
\]
\end{lemma}

\begin{theorem}{}
Определение меры Лебега корректно.
\end{theorem}

\begin{proofbox}{}
Пусть $\mu$ — мера на $\mathcal P$ ($\mathcal P$ — НЕ алгебра)

Доказательство счётной аддитивности:

Пусть \(P_k, P \subset \mathcal P\) и
\[
P=\bigsqcup_{k=1}^{\infty}P_k
\]

\[
P_k=\langle a_k,b_k\rangle,\qquad a_k,b_k\in\mathbb R^n,\qquad P=\langle a,b\rangle
\]

1) По лемме:
\[
\bigsqcup_{k=1}^{m}P_k\subset P
\quad\Longrightarrow\quad
\sum_{k=1}^{m}\mu P_k\leq \mu P
\]

При \(m\to\infty\):
\[
\sum_{k=1}^{\infty}\mu P_k\leq \mu P
\]

2) Для каждого \(P_k\): построим \(\widetilde P_k=\langle \widetilde a_k,\widetilde b_k\rangle\supset P_k\) (можем, потому что функция площади непрерывна) так, что
\[
\mu\widetilde P_k<\mu P_k+\frac{\varepsilon}{2^k}
\]

Имеем
\[
P\subset \overline P\subset \bigcup_{k=1}^{\infty}\widetilde P_k
\]
где \(P\) — компакт, \(\overline P\) — его замыкание, \(\widetilde P_k\) — открытые множества.

Выделяем конечное подпокрытие:
\[
\overline P\subset\bigcup_{k=1}^{m}\widetilde P_k
\]

По лемме (п. 2):
\[
\mu P=\mu\overline P\leq\sum_{k=1}^{m}\mu\widetilde P_k
\leq
\sum_{k=1}^{\infty}\mu P_k+\sum_{k=1}^{\infty}\frac{\varepsilon}{2^k}
\]

\[
\sum_{k=1}^{\infty}\mu P_k \leq \mu P \leq \sum_{k=1}^{\infty}\mu P_k+\varepsilon,
\qquad \text{верно при } \forall \varepsilon>0
\]

При \(\varepsilon\to 0\):
\[
\Rightarrow\quad
\sum_{k=1}^{\infty}\mu P_k\leq \mu P\leq\sum_{k=1}^{\infty}\mu P_k
\]

Отсюда
\[
\mu P=\sum_{k=1}^{\infty}\mu P_k
\]

Т.е. мы доказали счётную аддитивность.

\end{proofbox}

\subsection{Внутренняя и внешняя меры}

\begin{definition}
Пусть $A\subset\mathbb{R}^n$.

\textbf{Внешней мерой Лебега} называется

\[
\mu^*A=
\inf\left\{
\sum_{k=1}^{\infty}\mu(P_k):
A\subset\bigcup_{k=1}^{\infty}P_k,\
P_k\in\mathcal P
\right\}
\]

\begin{note}

Это ещё не мера.

\end{note}

\end{definition}

\begin{definition}
Пусть $A\subset\mathbb{R}^n$.

\textbf{Внутренней мерой Лебега} называется

\[
\mu_*A = \mu\Pi-\mu^*(\Pi-A)
\]

\begin{note}

1) Также можно определить внутреннюю меру таким образом:

\[
\mu^*A=
\sup\left\{
\sum_{k=1}^{\infty}\mu(P_k):
\bigcup_{k=1}^{\infty}P_k\subset A,\
P_k\in\mathcal P
\right\}
\]

2) Определение внутренней меры не зависит от бруса

3) Это тоже не мера

\end{note}

\end{definition}

\begin{definition}
Пусть $A\subset\mathbb{R}^n, A$ — ограничено.

$A$ называется измеримым по Лебегу, если

\[
\mu^*A=\mu_*A
\]

Тогда мера Лебега:

\[
\mu A=\mu^*A=\mu_*A
\]

\begin{note}

Брус измерим по Лебегу

\end{note}

\end{definition}

\begin{definition}
Пусть $A\subset\mathbb{R}^n, A$ — неограничено.

$A$ называется измеримым по Лебегу, если для любого $k\in\mathbb N$

\[
A\cap Q_k \forall k
\]

измеримо, где

\[
Q_k=[-k,k]^n \text{- кубы}
\]

Тогда

\[
\lambda A=\lim_{k\to\infty}\lambda(A\cap Q_k)
\]

\begin{note}

Предел $\exists$ в $[0; +\infty]$ по теореме Вейерштрасса, так как последовательность $\lambda(A\cap Q_k)$ возрастает

\end{note}

\end{definition}

\begin{theorem}{Каратеодори}
Семейство всех измеримых по Лебегу множеств в $\mathbb R^n$ является $\sigma$-алгеброй.

На этой $\sigma$-алгебре $\lambda$ является мерой.
\end{theorem}

\begin{lemma}{о структуре открытого множества}
Пусть $G\subset\mathbb R^n$ --- открытое.

Тогда существует счётное покрытие брусами:

\[
G=\bigcup_{i=1}^{\infty}P_i
\]

\end{lemma}

\begin{proofbox}{}

Пусть $x \in G, x$ — внутренняя $\Longrightarrow \exists B(x) \subset G \Longrightarrow \exists \: \text{брус} \: \Pi = [a;b) \subset B(x) \\ a, b \in \mathbb{Q}$

Получается счётное покрытие
\[
\bigcup_{i=1}^{\infty}P_i \supset G
\]

\[
\bigcup_{i=1}^{\infty}P_i = \bigsqcup_{i=1}^{\infty}(P_i\backslash P_{i-1}), P_0 = \emptyset, P_i\backslash P_{i — 1} = \bigsqcup \widetilde P_i
\]
Но $\bigcup_{i=1}^{\infty}P_i \subset G$

Таким образом

\[
\exists P_i (\widetilde P_i) : G = \bigcup_{i=1}^{\infty}P_i
\]

\end{proofbox}

\begin{note}
В $\mathbb R^n$ измеримы:

\begin{itemize}
\item брусы;
\item открытые множества;
\item замкнутые множества;
\item счётные объединения и пересечения измеримых множеств.
\end{itemize}
\end{note}

\begin{lemma}{об инвариантности относительно движения}
Мера Лебега не меняется относительно движения.

Пусть

\[
F:\mathbb R^n\to\mathbb R^n
\]

--- движение, а $A$ --- измеримое множество.

Тогда

\[
\mu(F(A))=\mu(A)
\]
\end{lemma}

\begin{example}
\textbf{\textit{Множество Витали}}
Пусть

\[
A=[0,1]\subset\mathbb R
\]

Рассмотрим отношение:

\[
x\sim y\Longleftrightarrow x-y\in\mathbb Q
\]

Пусть $A$ содержит по одному представителю каждого класса.

Тогда существуют $r_k\in[-1,1]\cap\mathbb Q$ такие, что

\[
[0, 1] \subset B = \bigcup_{k=1}^{\infty}(A+r_k) \subset [-1,2]
\]

Так как

\begin{flalign*}
\lambda B =\sum_{k=1}^\infty \lambda A =
\begin{cases}
0, & \text{если} $\lambda A = 0$,\\
$\infty$, & \text{если $\lambda A > 0$}
\end{cases}
\end{flalign*}

Но это невозможно. Следовательно, множество $A$ неизмеримо.
\end{example}

\newpage

\section{Регулярность меры}

\begin{theorem}{о регулярности меры}
Пусть $A \subset \mathbb{R}^n$, $A$ измеримо по Лебегу. Тогда:

\[
\lambda A = \inf\{\lambda G \mid G \text{ --- открыто},\ A \subset G\}
\]
(внешняя регулярность),

\[
\lambda A = \sup\{\lambda F \mid F \text{ --- замкнуто},\ F \subset A\}
\]
(внутренняя регулярность).
\end{theorem}

\begin{proofbox}{}
1) Докажем внешнюю регулярность.

Очевидно, что для любого открытого $G$, содержащего $A$,
\[
\lambda A \leq \lambda G
\]
по монотонности меры.

Пусть $\varepsilon>0$. Так как
\[
\lambda A
=
\inf\left\{
\sum_{k=1}^{\infty}\lambda P_k
\ \middle|\
A\subset\bigcup_{k=1}^{\infty}P_k,\ P_k\in\mathcal{P}
\right\}
\]
существуют брусы $P_k$ такие, что
\[
A\subset\bigcup_{k=1}^{\infty}P_k,
\qquad
\sum_{k=1}^{\infty}\lambda P_k<\lambda A+\varepsilon
\]

Для каждого $P_k$ выберем открытый брус $\widetilde P_k$ такой, что
\[
P_k\subset\widetilde P_k,
\qquad
\lambda\widetilde P_k<\lambda P_k+\frac{\varepsilon}{2^k}
\]

Обозначим
\[
G=\bigcup_{k=1}^{\infty}\widetilde P_k
\]
Тогда $G$ открыто и $A\subset G$. Кроме того,
\[
\lambda G
\leq
\sum_{k=1}^{\infty}\lambda\widetilde P_k
\leq
\sum_{k=1}^{\infty}\lambda P_k+\varepsilon
<
\lambda A+2\varepsilon
\]

Следовательно,
\[
\lambda A
=
\inf\{\lambda G\mid G\text{ --- открыто},\ A\subset G\}
\]

2) Докажем внутреннюю регулярность.

Сначала пусть $A$ ограничено. Тогда существует брус $\Pi$ такой, что
\[
A\subset\Pi
\]

Для любого $\varepsilon>0$ по внешней регулярности существует открытое множество $G$ такое, что
\[
\Pi\setminus A\subset G,
\qquad
\lambda G<\lambda(\Pi\setminus A)+\varepsilon
\]

Множество $\Pi\setminus G$ замкнуто и $\Pi\setminus G\subset A$

При этом
\[
\lambda(\Pi\setminus G)
=
\lambda\Pi-\lambda(G\cap\Pi)
>
\lambda\Pi-\lambda(\Pi\setminus A)-\varepsilon
=
\lambda A-\varepsilon
\]

Следовательно,
\[
\lambda A
=
\sup\{\lambda F\mid F\text{ --- замкнуто},\ F\subset A\}
\]

Теперь пусть $A$ неограничено. Тогда
\[
\lambda A
=
\lim_{k\to\infty}\lambda(A\cap Q_k)
=
\sup_k\lambda(A\cap Q_k)
\]

Для каждого $k$ множество $A\cap Q_k$ ограничено, поэтому
\[
\lambda(A\cap Q_k)
=
\sup\{\lambda F\mid F\subset A\cap Q_k,\ F\text{ --- замкнуто}\}
\]

Для каждого $k$ выберем замкнутое $F_k\subset A\cap Q_k$ так, что
\[
\lambda F_k>\lambda(A\cap Q_k)-\frac{1}{k}
\]

Можно считать, что
\[
F_k\subset F_{k+1}
\]

Тогда
\[
\lambda\left(\bigcup_{k=1}^{\infty}F_k\right)
=
\lim_{k\to\infty}\lambda F_k
\geq \lambda A
\]

Следовательно, для любого $\varepsilon>0$ найдётся замкнутое множество $F_m\subset A$ такое, что
\[
\lambda F_m>\lambda A-\varepsilon
\]

Значит,
\[
\lambda A
=
\sup\{\lambda F\mid F\text{ --- замкнуто},\ F\subset A\}
\]
\end{proofbox}

\begin{props}{}
Пусть $A$ измеримо по Лебегу. Тогда для любого $\varepsilon>0$ существуют открытое множество $G$ и замкнутое множество $F$ такие, что
\[
F\subset A\subset G
\]
\[
\lambda(G\setminus A)<\varepsilon,
\qquad
\lambda(A\setminus F)<\varepsilon
\]
\end{props}

\begin{note}
Вместо замкнутых множеств во внутренней регулярности можно брать компактные множества.
\end{note}

\newpage

\section{Мера Лебега--Стилтьеса}

\subsection{Опрделение}

В $\mathbb{R}$ в качестве полуинтервалов возьмём множества $[a,b)$, где
\[
a,b\in\mathbb{R},
\qquad
a\leq b
\]

Пусть
\[
f:\mathbb{R}\to\mathbb{R}
\]
причём $f$ не убывает.

\begin{definition}
На полуинтервалах $[a,b)$ зададим
\[
\mu_f[a,b)=f(b)-f(a)
\]
\end{definition}

Если $\mu_f$ является мерой, то она непрерывна. Пусть
\[
b_1<b_2<b_3<\ldots<b_n<b_{n+1}<\ldots,
\qquad
b_n\to b
\]

Тогда
\[
[a,b_n)\subset[a,b_{n+1})
\]
и
\[
\lim_{n\to\infty}\mu_f[a,b_n)
=
\mu_f\left(\bigcup_{n=1}^{\infty}[a,b_n)\right)
\]

Следовательно,
\[
\lim_{n\to\infty}\bigl(f(b_n)-f(a)\bigr)
=
f(b)-f(a)
\]
то есть
\[
\lim_{n\to\infty}f(b_n)=f(b)
\]

Таким образом,
\[
f(b-0)=f(b)
\]
то есть $f$ непрерывна слева.

\begin{note}
Если $f$ не убывает и непрерывна слева, то $\mu_f$ является мерой на системе полуинтервалов $[a,b)$.
\end{note}

\begin{note}
Если
\[
f(x)=x
\]
то
\[
\mu_f=\lambda
\]
то есть мера Лебега является частным случаем меры Лебега--Стилтьеса.
\end{note}

\subsection{Мера точки}

Имеем
\[
\bigcap_{k=1}^{\infty}\left[a,a+\frac{1}{k}\right)=\{a\}
\]

По непрерывности меры сверху
\[
\mu_f\{a\}
=
\lim_{k\to\infty}
\mu_f\left[a,a+\frac{1}{k}\right)
=
\lim_{k\to\infty}
\left(
f\left(a+\frac{1}{k}\right)-f(a)
\right)
\]

Следовательно,
\[
\mu_f\{a\}=f(a+0)-f(a)
\]

\begin{example}

\[
Пусть f(x)= \mathds{1}_{(0,+\infty)}(x)=
\begin{cases}
1, & x\in(0,+\infty),\\
0, & x\notin(0,+\infty).
\end{cases}
\]

Тогда мера $\mu_f$ сосредоточена в точке $0$.
\end{example}

Пусть $\lceil x\rceil=n$, то есть
\[
x\in(n-1,n]
\]

Для полуинтервалов получаем
\[
\mu_f[a,b)=f(b-0)-f(a+0)
\]
а для отрезков
\[
\mu_f[a,b]=f(b+0)-f(a)
\]

\newpage

\section{Измеримые функции}

\subsection{Определение}

Пусть $(X,\mathcal{A},\mu)$ --- измеримое пространство, где $X$ --- пространство, с которым мы работаем, $\mathcal{A}$ --- $\sigma$-алгебра измеримых подмножеств $X$, $\mu$ --- мера на $\mathcal{A}$.

\begin{definition}
\textbf{Измеримая функция} $f:X\to\overline{\mathbb{R}}$ --- функция, для которой все её лебеговы множества измеримы для любых a

\begin{align*}
\forall a\in\mathbb{R}\qquad
&E(f<a)=\{x\in X:f(x)<a\}\in\mathcal{A},\\
&E(f\leq a)=\{x\in X:f(x)\leq a\}\in\mathcal{A},\\
&E(f\geq a)=\{x\in X:f(x)\geq a\}\in\mathcal{A},\\
&E(f>a)=\{x\in X:f(x)>a\}\in\mathcal{A}.
\end{align*}

\end{definition}

\begin{theorem}{о лебеговых множествах}
Пусть $f:X\to\overline{\mathbb{R}}$. Тогда следующие условия эквивалентны:
\[
\forall a\in\mathbb{R}\quad E(f<a)\in\mathcal{A}
\]
\[
\forall a\in\mathbb{R}\quad E(f\leq a)\in\mathcal{A}
\]
\[
\forall a\in\mathbb{R}\quad E(f>a)\in\mathcal{A}
\]
\[
\forall a\in\mathbb{R}\quad E(f\geq a)\in\mathcal{A}
\]
\end{theorem}

\begin{proofbox}{}
Например,
\[
E(f\leq a)=\bigcap_{k=1}^{\infty}E\left(f<a+\frac1k\right)
\]
Так как все множества $E(f<a+\frac1k)$ измеримы, а $\mathcal{A}$ является $\sigma$-алгеброй, получаем измеримость пересечения.

Остальные случаи получаются аналогично с использованием дополнения множеств.
\end{proofbox}

\begin{note}
Если $f:\mathbb{R}\to\overline{\mathbb{R}}$ измерима, то множества E также измеримы.
\end{note}

\begin{example}

1) $f(x)=c$ на $E$ --- измерима, так как
\[
E(f<a)=
\begin{cases}
E,&a>c,\\
\varnothing,&a\leq c.
\end{cases}
\]

2) $f\in C(\mathbb{R})$ --- измерима, так как прообраз открытого множества при непрерывном отображении открыт.

3) Функция Дирихле:
\[
D(x)=
\begin{cases}
1,&x\in\mathbb{Q},\\
0,&x\in\mathbb{R}\setminus\mathbb{Q}
\end{cases}
\]
измерима.
\end{example}

Пусть $f:E\to\overline{\mathbb{R}}$ --- измеримая функция.

\begin{props}{измеримых функций}

1) Для любого $a\in\overline{\mathbb{R}}$
\[
f^{-1}(\{a\})
\]
измеримо.

2) Для любых $a,b\in\overline{\mathbb{R}}$
\[
f^{-1}(\langle a,b\rangle)
\]
измеримо.

3) Для любого открытого $G\subset\mathbb{R}$ множество
\[
f^{-1}(G)
\]
измеримо, а также замкнутого множества.

4) Если $A$ измеримо и $A\subset E$, то сужение
\[
f|_A
\]
измеримо.

5) Если $f|_{A_n}$ измерима $\forall n$, $n\in\mathbb{N}$, то $f$ измерима на
\[
\bigcup_{n=1}^{\infty}A_n
\]

6) Если $f$ является сужением некоторой функции на $X$, то она измерима.

\end{props}

\begin{theorem}{об измеримости композиции}
Пусть
\[
\varphi:\mathbb{R}^n\supset D\to\overline{\mathbb{R}},
\qquad
\varphi\in C(D)
\]
и
\[
f:\mathbb{R}\to D\subset\mathbb{R}^n
\]
где
\[
f=(f_1,f_2,\ldots,f_n)
\]
а все $f_i$ измеримы.

Тогда
\[
\varphi\circ f
\]
измерима.
\end{theorem}

\begin{proofbox}{}
Рассмотрим прообраз луча:
\[
(\varphi\circ f)^{-1}((-\infty,a))
=
f^{-1}\left(\varphi^{-1}((-\infty,a))\right)
\]

Так как $\varphi$ непрерывна, множество
\[
\varphi^{-1}((-\infty,a))
\]
является открытым в $D$.

Следовательно, его можно представить через счётное объединение:
\[
\bigcup_{i=1}^{\infty}[a_i,b_i)
\]
а прообразы этих множеств измеримы.

Значит, $(\varphi\circ f)$ измерима.
\end{proofbox}

\subsection{Операции над измеримыми функциями}

Пусть $f,g$ --- измеримые на $X$ функции.

\begin{props}

1) Функции
\[
f+g,\qquad f-g,\qquad fg,\qquad \frac fg\;(g\neq0)
\]
измеримы.

2) Функции
\[
\max(f,g),\qquad \min(f,g),\qquad |f|
\]
измеримы.

\end{props}

\begin{note}

1) $0 \: \cdot \: (\pm\infty) = 0$
2) $(\pm\infty) — (\pm\infty) = 0$
3) $(\mp\infty) + (\pm\infty) = 0$
4) $\frac{a}{\pm\infty} = 0 \;\; \forall a \in \overline{\mathbb{R}}$

\end{note}

\begin{theorem}{о предельном переходе}
Пусть
\[
f_n(x):X\to\overline{\mathbb{R}}
\]
измеримы. Тогда:

1)
\[
\inf f_n,\qquad \sup f_n
\]
измеримы.

2)
\[
\overline{\lim} f_n,\qquad \underline{\lim} f_n
\]
измеримы.

3) Если существует
\[
\lim f_n
\]
то предел измерим.
\end{theorem}

\begin{proofbox}{}
Например,
\[
\sup f_n
=
\lim_{n\to\infty}\sup_{k\leq n}f_k
\]

Так как конечные максимумы измеримы, то и предел измерим.

Аналогично рассматриваются $\inf$, $\limsup$ и $\liminf$.
\end{proofbox}

\newpage

\section{Простые функции}

\begin{definition}
\textbf{Простой функцией} называется измеримая функция
\[
f:X\to\mathbb{R}
\]
если она принимает конечное число значений.

Например:
\[
\operatorname{sign}(x),\qquad \cos(x)=c
\]

Пусть
\[
f(x)=
\begin{cases}
c_1,\ x\in A_1,\\
\vdots\\
c_n,\ x\in A_n,
\end{cases}
\]
где $A_i$ — измеримые дизъюнктные множества и
\[
\bigsqcup_{i=1}^{n}A_i=X
\]

Тогда $\{A_i\}$ — дизъюнктное разбиение множества $X$, и оно не единственно.
\end{definition}

Простую функцию можно записать в виде:
\[
f(x)=\sum_{k=1}^{n}c_k\cdot\chi_{A_k}(x)
\]
где
\[
\chi_{A}(x)=
\begin{cases}
1,\ x\in A,\\
0,\ x\notin A.
\end{cases}
\]

\begin{props}{простых функций}

1) Для любой простой функции существует такое дизъюнктное разбиение множества $X$, и оно не единственно.

2) Сумма, произведение, частное (знаменатель $\neq0$) конечного числа простых функций даёт снова простую функцию.

3) $\min$ и $\max$ конечного числа простых функций дают простую функцию.

4) Если $f$ — простая функция, то $f^{-1}((a,b))$, $a,b\in\overline{\mathbb{R}}$, измеримо.

\end{props}

\begin{theorem}{об аппроксимации измеримых функций простыми}
Пусть
\[
f:X\to[0,+\infty]
\]
измерима. Тогда существует возрастающая последовательность простых функций $f_n$, такая что
\[
\forall x\in X\qquad f_n(x)\to f(x)
\]
\end{theorem}

\begin{proofbox}{}
Пусть
\[
\Delta_i=
\left[
\frac{i}{2^n},
\frac{i+1}{2^n}
\right),
\qquad i=0,\ldots,n^2-1
\]
а также
\[
\Delta_{n^2}=[n,+\infty)
\]

Положим
\[
A_i=f^{-1}(\Delta_i)
\]

Тогда определим простую функцию
\[
g_n(x)=
\sum_{i=0}^{n^2}
\frac{i}{2^n}\chi_{A_i}(x)
\]

Но последовательность может быть невозрастающей

Тогда пусть $f_n(x) = \max\{g_1(x), g_2(x),...,g_n(x)\}$ — последовательность возрастает и стремится к f так как $0 \leq f-f_n \leq \frac{1}{n} \; \forall x \in X$

Следовательно, существует возрастающая последовательность простых функций, сходящаяся к $f$.
\end{proofbox}

\newpage

\section{Интеграл по мере}

Пусть $(X,\mathcal{A},\mu)$ — пространство с мерой

\begin{definition}
\textbf{Интегралом простой функции} называется число:
\[
\int_E f d\mu=
\sum_{k=1}^{n}c_k\mu(A_k)
\]
где
\[
f=\sum_{k=1}^{n}c_k\chi_{A_k}
\]
\end{definition}

\begin{definition}
Пусть
\[
f:X\supset E\to[0,+\infty]
\]
измерима.

Тогда
\[
\int_E f d\mu
:=
\sup\int_E g d\mu
\]
где $g$ — простая функция и
\[
0\leq g\leq f
\]
\end{definition}

\begin{definition}
Пусть
\[
f:X\supset E\to\overline{\mathbb{R}}
\]
измерима.

Положим
\[
f_+=\max\{f,0\}
\]
\[
f_-=\max\{-f,0\}
\]

Тогда
\[
f=f_+-f_-
\]

Интеграл определяется как
\[
\int_E f d\mu
:=
\int_E f_+ d\mu-
\int_E f_- d\mu
\]
\end{definition}

\begin{note}
Если оба интеграла
\[
\int_E f_+ d\mu,\qquad
\int_E f_- d\mu
\]
конечны, то $f$ называется суммируемой на $E$.

Если ровно один из интегралов равен $+\infty$, то
\[
\int_E f d\mu=\pm\infty
\]

Если оба интеграла бесконечны, то $f$ называется неинтегрируемой и
\[
\int_E f d\mu
\]
не существует.
\end{note}

\begin{example}
1) \[
f(x)=c=\operatorname{const}
\]

Если $c\geq0$, то
\[
\int_E c d\mu
=
c\cdot\mu(E)
\]

Если $c<0$, то
\[
\int_E c d\mu
=
(-c)\cdot\mu(E)
\]

\bigskip

2) Рассмотрим функцию Дирихле:
\[
D(x)=
\begin{cases}
1,\ x\in\mathbb{Q},\\
0,\ x\in\mathbb{R}\setminus\mathbb{Q}.
\end{cases}
\]

Пусть
\[
(\mathbb{R},\mathcal{A},\lambda)
\]
— пространство с мерой Лебега.

Тогда
\[
\int D(x) d\lambda
=
1\cdot\lambda(\mathbb{Q})
+
0\cdot\lambda(\mathbb{R}\setminus\mathbb{Q})
=0
\]
\end{example}

\begin{lemma}{о неучёте нулевых значений функции}
Пусть
\[
f:X\supset E\to\overline{\mathbb{R}}
\]
измерима, $E_1,\subset E$ измерима, причём
\[
f|_{E\backslash E_1}=0
\]

Тогда
\[
\int_{E}f d\mu=\int_{E_1}f_1d\mu
\]
Причём они существуют или не существуют одновременно
\end{lemma}

\begin{proofbox}{}
Если
\[
f=0
\]
то утверждение очевидно по определению.

Если
\[
f\geq0
\]
то по определению интеграла
\[
\int_E f d\mu
=
\sup\int_E g d\mu
\]
где $g$ — простая функция и $g\leq f$.

Тогда утверждение опять очевидно по 1 пункту

Для знакопеременной функции делим её на $f_+$ и $f_-$. Дальше действуем как в пункте 2
\end{proofbox}

\begin{corollary}
$\int_E fd\mu = \int_X f\:\cdot\:\chi_Ed\mu$
\end{corollary}

\begin{props}{}

1) Если $f=0$ на множестве, то интеграл по этому множеству равен нулю.

2) Если множества отличаются только на множество меры ноль, то интегралы совпадают.
\end{props}

\begin{theorem}{о монотонности по функции}
Пусть
\[
f,g:E\to\overline{\mathbb{R}}
\]
причём
\[
f\leq g\quad \forall x\in E
\]
и
\[
\int_E f d\mu, \int_E g d\mu \;\: \exists\; \text{в} \; \overline{\mathbb{R}}
\]

Тогда
\[
\int_E f d\mu\leq\int_E g d\mu
\]
\end{theorem}

\begin{proofbox}{}

1) Пусть $f$ и $g$ — простые функции, и
\[
\{A_k\}
\]
— общее допустимое разбиение.

Тогда
\[
f=\sum_{k=1}^{n}c_k\chi_{A_k},
\qquad
g=\sum_{k=1}^{n}d_k\chi_{A_k}
\]
где
\[
c_k\leq d_k
\]

Следовательно,
\[
\int_E f d\mu\leq\int_E g d\mu
\]

2) Пусть
\[
f,g\geq0
\]

Имеем:
\[
\{h\mid h\text{ простая},\ h\leq f\}
\subset
\{h\mid h\text{ простая},\ h\leq g\}
\]

Следовательно,
\[
\sup\{\ldots\}\leq\sup\{\ldots\}
\]
то есть
\[
\int_E f d\mu\leq\int_E g d\mu
\]

3) Пусть $f,g$ — знакопеременные

Имеем
\[
f_+\leq g_+,
\qquad
f_-\geq g_-
\]

Следовательно,
\[
\int_E f_+ d\mu-\int_E f_- d\mu
\leq
\int_E g_+ d\mu-\int_E g_- d\mu
\]

То есть
\[
\int_E f d\mu\leq\int_E g d\mu
\]

\end{proofbox}

\begin{corollary}
\textbf{Монотонность по множеству}

Пусть
\[
f:E\to[0,+\infty],
\qquad
E_1\subset E
\]

Тогда
\[
\int_{E_1}f d\mu
\leq
\int_Ef d\mu
\]

\end{corollary}

\begin{proofbox}{}

\[
\int_{E_1}f d\mu
=
\int_X f\chi_{E_1} d\mu
\]

\[
\int_{E}f d\mu
=
\int_X f\chi_{E} d\mu
\]

Но $E_1 \subset E$. Следовательно

\[
\int_{E_1}f d\mu
\leq
\int_Ef d\mu
\]

\end{proofbox}

\begin{corollary}
\textbf{Достаточное условие измеримости функции}

Пусть $f$ измерима и ограничена на $E$, причём
\[
\mu(E)<+\infty
\]

Тогда $f$ суммируема.
\end{corollary}

\begin{theorem}{Беппо Леви}
Пусть
\[
f_n:E\to[0,+\infty]
\]
— измеримые функции, такие что

\[
f_n(x)\leq f_{n+1}(x)
\quad
\forall x\in E
\]

и
\[
f_n(x)\to f(x)
\quad
\forall x\in E
\]

Тогда
\[
\int_E f_n d\mu
\longrightarrow
\int_E f d\mu
\]
\end{theorem}

\begin{proofbox}{}

1) a) f — измерима $\Longrightarrow \exists \int_E fd\mu$
б) Так как $f_n$ возрастают, то $\int f_n$ тоже возрастают и по теореме Вейерштрасса $\exists \lim\int_E f_nd\mu = L$ и
$\int f_nd\mu \leq \int fd\mu \Longrightarrow L \leq \int fd\mu$

2) Докажем, что $L \geq \int fd\mu$

Пусть $g$ — простая и $0 \leq g \leq f$, и

\[
g = \sum_{k=1}^m C_k \chi_{A_k}
\]

Пусть $\Theta \in (0, 1) $. Рассмотрим

\[
A_k \cap E_n = A_k \cap E(f_n\geq \Theta \cdot g)
\]

\[
E_n \subset E_{n+1}, \bigcup_{n=1}^{\infty}E_n = E
\]

Тогда

\[
\int_E f_nd\mu \geq \int_{E_n} f_nd\mu\geq \int_{E_n} \Theta\cdot gd\mu = \sum_{k=1}^n \Theta C_k\mu(A_k \cap E_n) \longrightarrow \sum_{k=1}^n \Theta C_k\mu A_k
\]

Следовательно

\[
L \geq \sum_{k=1}^n \Theta C_k\mu A_k = \int \Theta\cdot gd\mu
\]

Отсюда $L \geq \int fd\mu$

А так как $L \leq \int fd\mu$ и $L \geq \int fd\mu$, то $L = \int f_nd\mu = \int fd\mu$

\end{proofbox}

\begin{theorem}{о пренебрежимости множества меры ноль}
Пусть $f$ измерима на $E$ и
\[
E_1\subset E
\]
причём
\[
\mu(E\setminus E_1)=0
\]

Тогда
\[
\int_E f d\mu
=
\int_{E_1}f d\mu
\]
\end{theorem}

\begin{proofbox}{}
1) Если $f$ — простая функция,
\[
f=\sum c_k\chi_{A_k}
\]
то
\[
E=\bigcup A_k,
\qquad
E_1=\bigcup(A_k\cap E_1)
\]

Имеем:
\[
\mu(A_k)
=
\mu(A_k\cap E_1)+
\mu(A_k\cap(E\setminus E_1))
\]

Последний член равен нулю, поэтому
\[
\mu(A_k)=\mu(A_k\cap E_1)
\]

Следовательно,
\[
\int_E f d\mu
=
\int_{E_1}f d\mu
\]

2) Пусть
\[
f\geq0
\]

По теореме об аппроксимации существует последовательность простых функций $g_n$. По теореме Беппо Леви

\[
\int_Eg_n d\mu\to\int_Ef d\mu
\]
а также
\[
\int_{E_1}g_n d\mu\to\int_{E_1}f d\mu
\]

Так как для простых функций равенство уже доказано, получаем:
\[
\int_Ef d\mu
=
\int_{E_1}f d\mu
\]

3) Если f — знакопеременная, то рассматриваем $f_+$ и $f_-$ аналогично пункту 2

\end{proofbox}

\begin{definition}
Функции $f$ и $g$ называются \textbf{равными почти всюду}, если
\[
\mu\{x\in X:f(x)\neq g(x)\}=0
\]
Обозначается:
\[
f\sim g
\]
\end{definition}

\begin{theorem}{- Свойства интеграла}
Пусть $f,g$ измеримы на $E$ и
\[
f,g:E\to\overline{\mathbb{R}}
\]

Тогда:

1) $f \sim g \Longrightarrow \int_Efd\mu = \int_Egd\mu$



2) Линейность:
\[
\int_E(\alpha f+\beta g) d\mu
=
\alpha\int_Ef d\mu
+
\beta\int_Eg d\mu
\]
где
\[
\alpha,\beta\in\mathbb{R}
\]



3) Аддитивность:
если
\[
E_1\cup E_2=E,
\qquad
E_1\cap E_2=\emptyset
\]
то
\[
\int_Ef d\mu
=
\int_{E_1}f d\mu
+
\int_{E_2}f d\mu
\]



4) Если
\[
f>0,\qquad \mu(E)>0
\]
то
\[
\int_Ef d\mu>0
\]



5) Если
\[
\int_E|f| d\mu=0
\]
то
\[
f=0
\]
почти всюду на $E$.



6)$f$ суммируема тогда и только тогда, когда $|f|$ суммируема, причём
\[
\left|
\int_Ef d\mu
\right|
\leq
\int_E|f| d\mu
\]

\end{theorem}

\begin{proofbox}{свойства 1}
Так как
\[
f\sim g
\]
то существует множество
\[
A=\{x\in E:f(x)\neq g(x)\}
\]
для которого
\[
\mu(A)=0
\]

Следовательно,
\[
f=g
\]
на множестве
\[
E\setminus A
\]

Разобьём интеграл:

\[
\int_E f d\mu
=
\int_{E\setminus A}f d\mu
+
\int_A f d\mu
\]

Так как
\[
\mu(A)=0
\]
то

\[
\int_A f d\mu=0
\]

А на \(E\setminus A\):

\[
f=g
\]

поэтому:

\[
\int_{E\setminus A}f d\mu
=
\int_{E\setminus A}g d\mu
\]

Следовательно,

\[
\int_E f d\mu
=
\int_E g d\mu
\]
\end{proofbox}

\begin{proofbox}{свойства 2}
Пусть
\[
\int_E f d\mu,\qquad \int_E g d\mu
\]
существуют.

Если $f$ и $g$ — простые функции, то
\[
\alpha f+\beta g
\]
также является простой функцией (по свойству конечных сумм).

Пусть
\[
f,g\geq0
\]

Тогда по теореме об аппроксимации существуют последовательности простых функций:
\[
f_n\uparrow f,\qquad g_n\uparrow g
\]

Тогда
\[
\alpha f_n+\beta g_n\uparrow \alpha f+\beta g
\]

По теореме Леви:
\[
\int_E f_n d\mu\to\int_Ef d\mu
\]
\[
\int_E g_n d\mu\to\int_Eg d\mu
\]

Имеем:
\[
\int_E(\alpha f_n+\beta g_n) d\mu
=
\alpha\int_Ef_n d\mu+
\beta\int_Eg_n d\mu
\]

Переходя к пределу:
\[
\int_E(\alpha f+\beta g) d\mu
=
\alpha\int_Ef d\mu+
\beta\int_Eg d\mu
\]

Для произвольных функций раскладываем их на положительную и отрицательную части:
\[
f=f_+-f_-,
\qquad
g=g_+-g_-
\]
\end{proofbox}

\begin{proofbox}{свойства 3}
Пусть
\[
E_1\cup E_2=E,
\qquad
E_1\cap E_2=\emptyset
\]

Тогда
\[
\chi_E=\chi_{E_1}+\chi_{E_2}
\]

Следовательно:
\[
\int_E f d\mu
=
\int_Xf\chi_E d\mu
\]
и
\[
\chi_E=\chi_{E_1}+\chi_{E_2}
\]

Получаем:
\[
\int_Ef d\mu
=
\int_Xf\chi_{E_1} d\mu+
\int_Xf\chi_{E_2} d\mu
\]
то есть
\[
\int_Ef d\mu
=
\int_{E_1}f d\mu+
\int_{E_2}f d\mu
\]
\end{proofbox}

\begin{proofbox}{свойства 4}
Пусть
\[
f>0,\qquad \mu(E)>0
\]

Рассмотрим множества:
\[
E_n=\left\{x\in E:f(x)>\frac1n\right\}
\]

Тогда
\[
E=\bigcup_{n=1}^{\infty}E_n
\]

Если бы
\[
\mu(E_n)=0
\]
для всех $n$, то из счётной аддитивности:
\[
\mu(E)=0
\]
что противоречит условию.

Следовательно, существует $n$, для которого:
\[
\mu(E_n)>0
\]

Тогда:
\[
\int_Ef d\mu
\geq
\int_{E_n}f d\mu
\geq
\int_{E_n}\frac1n d\mu
=
\frac1n\mu(E_n)>0
\]
\end{proofbox}

\begin{proofbox}{свойства 5}
Пусть
\[
f\neq0
\]
на множестве
\[
A\subset E,
\qquad
\mu(A)>0
\]

Тогда:
\[
\int_E|f| d\mu
\geq
\int_A|f| d\mu>0
\]
по пункту 4.

Получили противоречие с условием:
\[
\int_E|f| d\mu=0
\]

Следовательно:
\[
f=0
\]
почти всюду.
\end{proofbox}

\begin{proofbox}{свойства 6}
Разложим функцию:

\[
f=f^+-f^-
\]

где

\[
|f|=f^++f^-
\]

Если

\[
\int_E |f| d\mu<\infty
\]

то

\[
\int_E f^+ d\mu
+
\int_E f^- d\mu
<\infty
\]

Следовательно,

\[
\int_E f^+ d\mu<\infty
\]

и

\[
\int_E f^- d\mu<\infty
\]

Поэтому существует:

\[
\int_E f d\mu
=
\int_E f^+ d\mu
-
\int_E f^- d\mu
\]

То есть $f$ суммируема.

Обратно, если $f$ суммируема, то:

\[
\int_E f^+ d\mu<\infty
\]

\[
\int_E f^- d\mu<\infty
\]

Следовательно:

\[
\int_E|f| d\mu
=
\int_E(f^++f^-) d\mu
<\infty
\]

Значит, $|f|$ суммируема.

Теперь докажем неравенство.

Так как:

\[
-f^-\leq f\leq f^+
\]

то:

\[
-\int_E f^-d\mu
\leq
\int_E f d\mu
\leq
\int_E f^+d\mu
\]

Следовательно:

\[
\left|\int_E f d\mu\right|
\leq
\max
\left(
\int_Ef^+d\mu,
\int_Ef^-d\mu
\right)
\]

А значит:

\[
\left|\int_Ef d\mu\right|
\leq
\int_E(f^++f^-) d\mu
\]

То есть:

\[
\left|\int_Ef d\mu\right|
\leq
\int_E|f| d\mu
\]
\end{proofbox}

\begin{theorem}{Неравенство Чебышёва}
Пусть $f$ измерима на $E$, $t>0$.

Тогда:
\[
\mu\{x\in E:|f(x)|\geq t\}
\leq
\frac1t\int_E|f| d\mu
\]
\end{theorem}

\begin{proofbox}{}
Обозначим:
\[
E_t=\{x\in E:|f(x)|\geq t\}
\]

Тогда:
\[
\int_E|f| d\mu
\geq
\int_{E_t}|f| d\mu
\]

Так как на $E_t$:
\[
|f|\geq t
\]
то:
\[
\int_{E_t}|f| d\mu
\geq
\int_{E_t}t d\mu
=
t\mu(E_t)
\]

Следовательно:
\[
t\mu(E_t)\leq\int_E|f| d\mu
\]

Делим на $t$:
\[
\mu(E_t)
\leq
\frac1t\int_E|f| d\mu
\]
\end{proofbox}

\begin{corollary}
Если $f$ суммируема, то $f$ конечно почти везде
\end{corollary}

\newpage

\section{Сравнение интегралов Римана и Лебега}

\begin{theorem}{Критерий Лебега интегрируемости по Риману}
Пусть
\[
f:[a,b]\to\mathbb{R}
\]
ограничена.

Тогда:
\[
f\in R[a,b]
\Longleftrightarrow
\lambda\{x\in[a,b]:f\text{ имеет разрыв в }x\}=0
\]
\end{theorem}

\begin{theorem}{о связи интегралов Лебега и Римана}
Если
\[
f\in R[a,b]
\]
то существует интеграл Лебега и:
\[
\int_{[a,b]}f d\lambda
=
\int_a^bf(x) dx
\]
\end{theorem}

\begin{proofbox}{}
Пусть
\[
\tau_n=\{x_i\}
\]
— разбиение отрезка на $2^n$ частей.

Обозначим:
\[
m_i=\inf_{x\in[x_{i-1},x_i]}f(x)
\]
\[
M_i=\sup_{x\in[x_{i-1},x_i]}f(x)
\]

Тогда нижняя и верхняя суммы:
\[
s_n=\sum_{i=1}^{n}m_i\Delta x_i
\]
\[
S_n=\sum_{i=1}^{n}M_i\Delta x_i
\]

Определим простые функции:
\[
\underline f_n(x)=
\sum_{i=1}^{n}
m_i\chi_{[x_{i-1},x_i]}(x)
\]

\[
\overline f_n(x)=
\sum_{i=1}^{n}
M_i\chi_{[x_{i-1},x_i]}(x)
\]

При измельчении разбиения:
\[
\underline f_n\uparrow f,
\qquad
\overline f_n\downarrow f
\]
и ограничены (сверху и снизу соответственно). Следовательно

\[
\exists \underline f(x) = \lim \underline f_n
\text{и}
\exists \overline{f(x)} = \lim \overline{f_n}
\]

Тогда по теореме Беппо Леви

\[
\int_{[a, b]}f_nd\lambda\longrightarrow\int\underline fd\lambda
\]

Для $\overline{f_n}: g_n = \overline{f_1} — \overline{f_n} \longrightarrow \overline{f_1} — \overline{f} \Longrightarrow \int \overline{f_n} \longrightarrow \int \overline{f}$

Но
\[
\int \underline{f_n} = s_n, \int \overline{f_n} = S_n
\]
\[
S_n - s_n \longrightarrow 0 \Longrightarrow \int(\overline{f_n} - \underline{f_n}) \longrightarrow \int(\overline{f} - \underline{f}) = 0 \Longrightarrow
\]
\[
\Longrightarrow \overline{f} = \underline{f} \:\ \text{п.в.}
= f  \:\ \text{п.в.} \Longrightarrow \int fd\lambda
\]

\end{proofbox}

\newpage

\section{Интеграл как функция множества}

Пусть $f\colon X\to\overline{\mathbb{R}}$, $f$ --- фиксированная, $f$ --- измеримая функция.

\[
\varphi(A)=\int_A f\,d\mu,
\qquad
A\subset X,\quad A\text{ --- измеримое}
\]

\begin{theorem}{о счётной аддитивности интеграла}

Пусть $f$ суммируема на $E$, то есть $f\in L(E)$,
\[
E=\bigsqcup_{k=1}^{\infty}E_k
\]
Тогда
\[
\int_E f\,d\mu
=
\sum_{k=1}^{\infty}\int_{E_k}f\,d\mu
\]

\end{theorem}

\begin{proofbox}{}

1) Пусть $f\geq 0$.

\[
\chi_E=\sum_{k=1}^{\infty}\chi_{E_k},
\qquad
f\chi_E=\sum_{k=1}^{\infty}f\chi_{E_k}
\]

Положим
\[
S_n(x)=\sum_{k=1}^{n}f(x)\chi_{E_k}(x)\geq 0
\]
Тогда
\[
S_n\uparrow,
\qquad
S_n\longrightarrow f\chi_E
\]

По теореме Леви
\[
\int_E S_n\,d\mu\longrightarrow\int_E f\,d\mu
\]
но
\[
\int_E S_n\,d\mu
=
\sum_{k=1}^{n}\int_{E_k}f\,d\mu
\]

2) Если $f$ произвольная, то
\[
f=f_+-f_-
\]

\end{proofbox}

\begin{corollary}

Если $f\colon X\to[0,+\infty]$, $f\in L(X)$, то
\[
\varphi(A)=\int_A f\,d\mu
\]
--- мера, а $f$ --- плотность меры.

\end{corollary}

\begin{corollary}

\[
\int_A f\,d\mu
\]
обладает свойством абсолютной непрерывности меры, а именно: непрерывностью сверху и снизу.

\end{corollary}

\begin{corollary}

Если $f\in R_{\mathrm{loc}}[a,+\infty)$, $f\geq0$, то
\[
\int_a^{+\infty}f\,dx
\]
сходится. Тогда
\[
\int_{[a,+\infty)}f\,d\lambda
=
\int_a^{+\infty}f\,dx
\]

\end{corollary}

\begin{proofbox}{}

\[
\int_a^{+\infty}f\,dx
=
\lim_{n\to+\infty}\int_a^n f\,dx
=
\lim_{n\to+\infty}\int_{[a,n]}f\,d\lambda
=
\int_{[a,+\infty)}f\,d\lambda
\]

\end{proofbox}

\begin{corollary}

Если
\[
\int_a^{+\infty}f\,dx
\]
сходится абсолютно, то
\[
\int_{[a,+\infty)}f\,d\lambda
=
\int_a^{+\infty}f\,dx
\]

Если несобственный интеграл сходится условно, то один из $\int f_{\pm}$ конечен, другой --- бесконечен.

\end{corollary}

\begin{theorem}{- Абсолютная непрерывность интеграла}

Пусть $f$ --- суммируемая функция, то есть $f\in L(X)$. Тогда для любого $\varepsilon>0$ существует $\delta>0$ такое, что
\[
\mu(A)<\delta
\quad\Longrightarrow\quad
\int_A|f|\,d\mu<\varepsilon
\]

\end{theorem}

\begin{proofbox}{}

Так как
\[
\int_E|f|\,d\mu
=
\sup\left\{
\int_E g\,d\mu:
0\leq g\leq|f|,
\quad g\text{ --- простая}
\right\}
\]
то для любого $\varepsilon>0$ существует $g_\varepsilon$ такая, что
\[
\int_E g_\varepsilon\,d\mu
>
\int_E|f|\,d\mu-\frac{\varepsilon}{2}
\]

Пусть
\[
|g_\varepsilon|\leq C_\varepsilon
\]
Возьмём
\[
\delta=\frac{\varepsilon}{2C_\varepsilon}
\]

Если $\mu(A)<\delta$, то
\[
\begin{aligned}
\int_A|f|\,d\mu
&=
\int_A\bigl(|f|-g_\varepsilon\bigr)\,d\mu
+
\int_Ag_\varepsilon\,d\mu
\\
&\leq
\frac{\varepsilon}{2}
+
C_\varepsilon\frac{\varepsilon}{2C_\varepsilon}
=
\varepsilon.
\end{aligned}
\]

\end{proofbox}

\newpage

\section{Функциональные последовательности и ряды}

Пусть $X$ --- метрическое, нормированное или вообще произвольное пространство (в $\mathbb{R}^m\text{ или }\mathbb{R})$.

Пусть
\[
f_n(x)\colon X\supset E\to\mathbb{R}
\]
--- \textbf{функциональная последовательность}:
\[
f_1(x),\,f_2(x),\,\ldots
\]

\begin{definition}{Множество сходимости функциональной последовательности}

\[
D=\left\{x\in E:\exists\lim_{n\to\infty}f_n(x)\in\mathbb{R}\right\}
\]

На $D$ определена функция
\[
f(x)=\lim_{n\to\infty}f_n(x)
\]

$f(x)$ --- предельная функция.

\end{definition}

\begin{example}

Пусть
\[
f_n(x)=x^n:\quad 1,x,x^2,x^3,\ldots,
\qquad E=\mathbb{R},\qquad D=(-1,1]
\]
и
\[
\lim_{n\to\infty}x^n=
\begin{cases}
0, & x\in(-1,1),\\
1, & x=1.
\end{cases}
\]

\end{example}

\begin{definition}{Множество сходимости функционального ряда}

Пусть функциональный ряд
\[
\sum_{k=1}^{\infty}f_k(x)
\]
где
\[
f_k(x)\colon X\supset E\to\mathbb{R}
\]

\[
D=
\left\{
x\in E:
\exists\lim_{n\to\infty}\sum_{k=1}^{n}f_k(x)\in\mathbb{R}
\right\}
\]

\end{definition}

\begin{example}

1)
\[
\sum_{n=0}^{\infty}x^n
\]
сходится при $x\in(-1,1)$:
\[
\sum_{n=0}^{\infty}x^n=\frac{1}{1-x}
\]

2)
\[
\sum_{n=1}^{\infty}\frac{x^n}{n}
\]

При $|x|>1$ ряд расходится, так как
\[
\frac{x^n}{n}\not\to0
\]

При $x\in(-1,1)$:
\[
\left|\frac{x^n}{n}\right|\leq |x|^n
\]
поэтому ряд сходится.

При $x=1$:
\[
\sum_{n=1}^{\infty}\frac{1}{n}
\]
расходится, а при $x=-1$:
\[
\sum_{n=1}^{\infty}\frac{(-1)^n}{n}
\]
сходится.

Следовательно,
\[
D=[-1,1)
\]

\end{example}

\newpage

\section{Равномерная сходимость}

\subsection{Определение}

\begin{definition}{Равномерная сходимость}

Пусть
\[
f_n(x)\colon E\to\mathbb{R}
\]

Последовательность $f_n(x)$ сходится равномерно на $D\subset E$ к $f(x)$,
\[
f_n(x)\overset{D}{\rightrightarrows}f(x) \Longleftrightarrow
\]
\[
\Longleftrightarrow\forall\varepsilon>0\ \exists n_0\ 
\forall n\geq n_0\ 
\forall x\in D:
\quad
|f_n(x)-f(x)|<\varepsilon
\]

\end{definition}

\begin{definition}
Поточечная сходимость на $D$:
\[
\forall x\in D\ \forall\varepsilon>0\ \exists n_0(\varepsilon,x)\
\forall n\geq n_0:
\quad
|f_n(x)-f(x)|<\varepsilon
\]
\end{definition}

\begin{example}

В равномерной сходимости все функции должны попасть внутрь трубки.

Для поточечной сходимости для каждого $x$ разный номер $n_0$.

\end{example}

\begin{note}

Если
\[
f_n\overset{D}{\rightrightarrows}f
\]
то
\[
f_n\overset{D}{\longrightarrow}f
\]
поточечно.

\end{note}

\begin{example}

Пусть
\[
f_n(x)=x^n
\longrightarrow
\begin{cases}
0, & x\in(-1,1),\\
1, & x=1.
\end{cases}
\]

1) На
\[
D_1=(-1,1]
\]
поточечная сходимость.

2) Пусть
\[
D_2=[-a,a]\subset(-1,1)
\]
На $D_2$
\[
|x^n-0|\leq a^n<\varepsilon
\]

Можно взять
\[
n_0=\left[\log_a\varepsilon\right]+1
\]
то есть $n_0$ подходит для всех $x$. Следовательно,
\[
x^n\xrightarrow[n\to\infty]{D_2}0
\]

Покажем, что на
\[
D_3=(-1,1)
\]
сходимость к нулю неравномерная.

Отрицание равномерной сходимости:
\[
\exists\varepsilon>0\ \forall n_0\ 
\exists n\geq n_0\ 
\exists x\in D:
\quad
|f_n(x)-f(x)|\geq\varepsilon
\]

Возьмём
\[
x_n=1-\frac1n
\]
Тогда
\[
f_n(x_n)=\left(1-\frac1n\right)^n
\longrightarrow\frac1e>\frac{1}{2e}=\varepsilon
\]

Значит, последовательность $x^n$ сходится к нулю неравномерно на
\[
D=(-1,1)
\]

\end{example}

Введём норму для $f\colon E\to\mathbb{R}$:
\[
\|f\|_E:=\sup_{x\in E}|f(x)|
\]
--- \textbf{суп-норма (равномерная, чебышёвская).}

\begin{note}

\[
f_n\overset{D}{\rightrightarrows}f
\quad\Longleftrightarrow\quad
\|f_n-f\|_D\longrightarrow0,
\qquad n\to\infty
\]

\end{note}

\begin{example}

Пусть
\[
f_n(x)=x^n,
\qquad
D_2=[-a,a]\subset(-1,1),
\qquad
D_3=(-1,1)
\]

Тогда
\[
\|x^n\|_{D_2}
=
\sup_{x\in[-a,a]}|x^n|
=
a^n\longrightarrow0
\]
а
\[
\|x^n\|_{D_3}
=
\sup_{x\in(-1,1)}|x^n|
=
1\longrightarrow1\neq0
\]

\end{example}

\begin{definition}

\textbf{Равномерная сходимость функционального ряда}

Пусть
\[
f_n\colon E\to\mathbb{R}
\]

Ряд
\[
\sum_{n=1}^{\infty}f_n(x)
\]
сходится равномерно к $S(x)$, если
\[
S_n(x)=\sum_{k=1}^{n}f_k(x)\xrightarrow{D}S(x)
\]

\end{definition}

\begin{note}

\[
\sum_{k=1}^{\infty}f_k(x)\overset{D}{\rightrightarrows}S(x)
\quad\Longleftrightarrow\quad
\|S_n-S\|_D\longrightarrow0
\quad\Longleftrightarrow\quad
\|R_n\|_D\longrightarrow0
\quad\Longleftrightarrow\quad
R_n(x)\overset{D}{\rightrightarrows}0
\]

\end{note}

\begin{theorem}{Критерий Коши равномерной сходимости функциональной последовательности}

\[
f_n(x)\overset{D}{\rightrightarrows}f(x)
\quad\Longleftrightarrow\quad
\forall\varepsilon>0\ \exists n_0\
\forall n\geq n_0\
\forall p\in\mathbb{N}\
\forall x\in D:
\quad
|f_{n+p}(x)-f_n(x)|<\varepsilon
\]

\end{theorem}

\begin{proofbox}{}

\textbf{Необходимость.}

Пусть
\[
f_n(x)\overset{D}{\rightrightarrows}f(x)
\]
Тогда
\[
\forall\varepsilon>0\ \exists n_0\
\forall n\geq n_0\
\forall x\in D:
\quad
|f_n(x)-f(x)|<\frac{\varepsilon}{2}
\]
Следовательно,
\[
|f_{n+p}(x)-f_n(x)|
\leq
|f_{n+p}(x)-f(x)|
+
|f(x)-f_n(x)|
<\varepsilon
\]

\textbf{Достаточность.}

Пусть
\[
\forall\varepsilon>0\ \exists n_0\
\forall n\geq n_0\
\forall p\in\mathbb{N}\
\forall x\in D:
\quad
|f_{n+p}(x)-f_n(x)|<\varepsilon
\]

Тогда для каждого $x\in D$ последовательность $f_n(x)$ фундаментальна (выполнен критерий Коши для последовательности), следовательно, сходится:
\[
\exists f(x)=\lim_{n\to\infty}f_n(x)
\]

Из
\[
|f_{n+p}(x)-f_n(x)|<\varepsilon
\]
при $p\to\infty$ получаем
\[
|f(x)-f_n(x)|\leq\varepsilon
\]
а это определение равномерной сходимости. Следовательно,
\[
f_n\rightrightarrows f
\]

\end{proofbox}

\begin{theorem}{Критерий Коши равномерной сходимости функционального ряда}

Ряд
\[
\sum_{k=1}^{\infty}f_k(x)
\]
сходится равномерно на $D$ тогда и только тогда, когда
\[
\forall\varepsilon>0\ \exists n_0\
\forall n\geq n_0\
\forall p\in\mathbb{N}\
\forall x\in D:
\quad
\left|
\sum_{k=n+1}^{n+p}f_k(x)
\right|<\varepsilon
\]

\end{theorem}

\begin{proofbox}{}

Очевидно, так как
\[
S_{n+p}(x)-S_n(x)
=
\sum_{k=n+1}^{n+p}f_k(x)
\]

\end{proofbox}

\begin{corollary}

\textbf{Необходимое условие равномерной сходимости ряда}

Если
\[
\sum_{k=1}^{\infty}f_k(x)
\]
сходится равномерно на $D$, то
\[
f_k(x)\overset{D}{\rightrightarrows}0
\]

\end{corollary}

\subsection{Простые свойства равномерно сходящихся функциональных последовательностей}

Пусть
\[
f_n,g_n\colon D\to\mathbb{R}
\]

\begin{props}

1) Если
\[
f_n\overset{D}{\rightrightarrows}f,
\qquad
g_n\overset{D}{\rightrightarrows}g
\]
то для любых $\alpha,\beta\in\mathbb{R}$
\[
\alpha f_n+\beta g_n\overset{D}{\rightrightarrows}\alpha f+\beta g
\]

2) Если $D_1\subset D$ и
\[
f_n\overset{D}{\rightrightarrows}f \text{, то}\ f_n\overset{D_1}{\rightrightarrows}f
\]

3) Если
\[
f_n\overset{D}{\rightrightarrows}f \text{ и } g \text{ ограничена на } D\text{, то } f_ng\overset{D}{\rightrightarrows}fg
\]

\end{props}

\begin{proofbox}{свойства 1}

Пусть
\[
f_n\overset{D}{\rightrightarrows}f,
\qquad
g_n\overset{D}{\rightrightarrows}g
\]

Для любых $\alpha,\beta\in\mathbb{R}$ имеем
\[
\begin{aligned}
\left|
(\alpha f_n(x)+\beta g_n(x))
-
(\alpha f(x)+\beta g(x))
\right|
&\leq
|\alpha|\,|f_n(x)-f(x)|
+
|\beta|\,|g_n(x)-g(x)|.
\end{aligned}
\]

Так как $f_n\overset{D}{\rightrightarrows}f$ и $g_n\overset{D}{\rightrightarrows}g$, правая часть стремится к нулю равномерно по $x\in D$. Следовательно,
\[
\alpha f_n+\beta g_n
\overset{D}{\rightrightarrows}
\alpha f+\beta g
\]

\end{proofbox}

\begin{proofbox}{свойства 2}

Пусть $D_1\subset D$ и
\[
f_n\overset{D}{\rightrightarrows}f
\]

Тогда
\[
\forall\varepsilon>0\ \exists n_0\
\forall n\geq n_0\
\forall x\in D:
\quad
|f_n(x)-f(x)|<\varepsilon
\]

Так как $D_1\subset D$, то тем более
\[
\forall\varepsilon>0\ \exists n_0\
\forall n\geq n_0\
\forall x\in D_1:
\quad
|f_n(x)-f(x)|<\varepsilon
\]

Следовательно,
\[
f_n\overset{D_1}{\rightrightarrows}f
\]

\end{proofbox}

\begin{proofbox}{свойства 3}

Пусть
\[
f_n\overset{D}{\rightrightarrows}f
\]
и функция $g$ ограничена на $D$, то есть
\[
|g(x)|\leq C
\qquad
\forall x\in D
\]

Тогда
\[
|f_n(x)g(x)-f(x)g(x)|
=
|g(x)|\,|f_n(x)-f(x)|
\leq
C|f_n(x)-f(x)|
\]

Так как $f_n\overset{D}{\rightrightarrows}f$, для любого $\varepsilon>0$ начиная с некоторого номера $n_0$
\[
|f_n(x)-f(x)|<\frac{\varepsilon}{C}
\qquad
\forall x\in D
\]

Следовательно,
\[
|f_n(x)g(x)-f(x)g(x)|<\varepsilon
\qquad
\forall x\in D
\]
то есть
\[
f_ng\overset{D}{\rightrightarrows}fg
\]

\end{proofbox}

\begin{note}

Для рядов эти свойства тоже верны.

\end{note}

\begin{example}

Рассмотрим
\[
\sum_{n=1}^{\infty}x^n,
\qquad
D=(-1,1)
\]
Так как
\[
f_n(x)=x^n\overset{D}{\rightrightarrows}0
\]
неравномерно, необходимое условие равномерной сходимости не выполнено.

Следовательно,
\[
\sum_{n=1}^{\infty}x^n
\]
сходится неравномерно.

\end{example}

\newpage

\section{Признаки равномерной сходимости ряда}

\begin{theorem}{Признак Вейерштрасса}

Пусть
\[
f_n(x)\colon X\to\mathbb{R},
\qquad
D\subset X
\]
и для всех $x\in D: |f_n(x)|\leq a_n \text{ и } \sum_{n=1}^{\infty}a_n$ сходится, то ряд
\[
\sum_{n=1}^{\infty}f_n(x)
\]
сходится равномерно на $D$ и абсолютно.

$a_n$ --- \textbf{мажоранта}

\end{theorem}

\begin{proofbox}{}

По критерию Коши
\[
\sum f_n(x)
\]
сходится равномерно тогда и только тогда, когда
\[
\forall\varepsilon>0\ \exists n_0\
\forall n\geq n_0\
\forall p\in\mathbb{N}\
\forall x\in D:
\quad
\left|
\sum_{k=n+1}^{n+p}f_k(x)
\right|<\varepsilon
\]

Имеем
\[
\left|
\sum_{k=n+1}^{n+p}f_k(x)
\right|
\leq
\sum_{k=n+1}^{n+p}|f_k(x)|
\leq
\sum_{k=n+1}^{n+p}a_k
<\varepsilon
\]
по критерию Коши для ряда $\sum a_k$.

То есть выполнен критерий Коши равномерной сходимости ряда.

\end{proofbox}

\begin{example}

Ряд
\[
\sum_{n=1}^{\infty}x^n
\]
сходится на $(-1,1)$.

Пусть
\[
D=[-a,a]\subset(-1,1)
\]
Тогда
\[
|x^n|\leq a^n
\]

Следовательно, по признаку Вейерштрасса ряд сходится равномерно на $D$.

\end{example}

\begin{theorem}{Признаки Дирихле--Абеля равномерной сходимости ряда}

Рассмотрим ряд
\[
\sum_{n=1}^{\infty}f_n(x)g_n(x),
\qquad
f_n,g_n\colon X\to\mathbb{R},
\qquad
D\subset X
\]

Для равномерной сходимости на $D$ ряда достаточно выполнения одного из следующих наборов условий.

\textbf{Признак Дирихле:}

1) Частичные суммы
\[
A_n(x)=\sum_{k=1}^{n}f_k(x)
\]
равномерно ограничены, то есть существует $C$ такое, что
\[
|A_n(x)|\leq C
\qquad
\forall n,\quad\forall x\in D
\]

2) $g_n(x)$ монотонна для каждого $x\in D$ и
\[
g_n(x)\xrightarrow[n\to\infty]{D}0
\]

\textbf{Признак Абеля:}

1) Ряд
\[
\sum_{k=1}^{\infty}f_k(x)
\]
сходится равномерно на $D$.

2) $g_n(x)$ монотонна для каждого $x\in D$ и равномерно ограничена:
\[
|g_n(x)|\leq C
\]

\end{theorem}

\begin{proofbox}{}

\textbf{Признак Дирихле.}

Положим
\[
A_n=\sum_{k=1}^{n}a_k
\]
По преобразованию Абеля
\[
\sum_{k=1}^{n}a_kb_k
=
A_nb_n+
\sum_{k=1}^{n-1}A_k(b_k-b_{k+1})
\]

Поэтому
\[
\begin{aligned}
\left|
\sum_{k=n+1}^{n+p}f_k g_k
\right|
&=
\left|
\sum_{k=1}^{n+p}f_k g_k
-
\sum_{k=1}^{n}f_k g_k
\right|
\\
&=
\left|
A_{n+p}g_{n+p}
+
\sum_{k=n}^{n+p-1}A_k(g_k-g_{k+1})
-
A_ng_n
-
\sum_{k=1}^{n-1}A_k(g_k-g_{k+1})
\right|
\\
&\leq
|A_{n+p}|\,|g_{n+p}|
+
|A_n|\,|g_n|
+
\sum_{k=n}^{n+p-1}|A_k|\,g_k-g_{k+1}|
\\
&\leq
C\left(
|g_{n+p}|+|g_n|
+
\sum_{k=n}^{n+p-1}|g_k-g_{k+1}|
\right).
\end{aligned}
\]

Так как $g_n$ монотонна, то
\[
\sum_{k=n}^{n+p-1}|g_k-g_{k+1}|
=
|g_n-g_{n+p}|
\]
Следовательно,
\[
\left|
\sum_{k=n+1}^{n+p}f_k g_k
\right|
\leq
C\bigl(|g_{n+p}|+|g_n|+|g_n-g_{n+p}|\bigr)
\leq
2C\bigl(|g_{n+p}|+|g_n|\bigr)
\]

Так как
\[
g_n\overset{D}{\rightrightarrows}0
\]
то для всех достаточно больших $n$
\[
|g_n|<\frac{\varepsilon}{4C}
\]
равномерно по $x\in D$. Поэтому
\[
\left|
\sum_{k=n+1}^{n+p}f_k(x)g_k(x)
\right|<\varepsilon
\]

По критерию Коши ряд сходится равномерно на $D$.

\bigskip

\textbf{Признак Абеля.}

Так как ряд
\[
\sum_{k=1}^{\infty}f_k(x)
\]
сходится равномерно, то его остаток
\[
R_n(x)=\sum_{k=1}^{\infty}f_k(x) - A_n(x) = S(x) - A_n(x)
\]
удовлетворяет
\[
R_n(x)\xrightarrow[n\to\infty]{D}0
\]

Положим
\[
A_n=S-R_n
\]
где $S$ --- сумма ряда. Применяя преобразование Абеля к хвосту ряда, получаем
\[
\begin{aligned}
\left|
\sum_{k=n+1}^{n+p}f_k g_k
\right|
&=
|S - R_{n+p}|\,|g_{n+p}|
+
|S - R_n|\,|g_n|
+
\sum_{k=n}^{n+p-1}|S - R_k|\,|g_k-g_{k+1}|
\\
&=
\left|
S(g_{n+p}-g_n+\sum_{k=n}^{n+p-1}g_k-g_{k+1} + ...)
\right|
\\
&\leq
|R_{n+p}| \cdot |g_{n+p}| + |R_n| \cdot |g_n| + \sum_{k=n}^{n+p-1}R_k(g_k-g_{k+1})
\\
&\leq
C\frac{\varepsilon}{4C}
+
\frac{\varepsilon}{4C}C
+
2C\frac{\varepsilon}{4C}
=
\varepsilon.
\end{aligned}
\]

По критерию Коши ряд сходится равномерно на $D$.

\end{proofbox}

\begin{example}

Рассмотрим ряд
\[
\sum_{n=1}^{\infty}\frac{\sin(nx)}{n}
\]

Положим
\[
A_n(x)=\sum_{k=1}^{n}\sin(kx)
\]
Тогда
\[
A_n(x)
=
\frac{\sin\frac{nx}{2}\,
\sin\frac{(n+1)x}{2}}
{\sin\frac{x}{2}}
\]

Если
\[
x\in(\delta,2\pi-\delta),
\qquad
\delta\in(0,\pi)
\]
то
\[
\left|\sin\frac{x}{2}\right|
>
\sin\frac{\delta}{2}
\]
Поэтому
\[
|A_n(x)|
<
\frac{1}{\sin\frac{\delta}{2}}
\]

Кроме того,
\[
g_n(x)=\frac1n\downarrow0
\]
причём сходимость к нулю равномерная на $D$.

Следовательно, по признаку Дирихле ряд
\[
\sum_{n=1}^{\infty}\frac{\sin(nx)}{n}
\]
сходится равномерно на $[\delta,2\pi-\delta]$.

\end{example}

\begin{lemma}

Пусть
\[
f_k(x)\in C[a,b],
\qquad
f_k(x)\xrightarrow[k\to\infty]{[a,b]}f(x)
\]

Тогда $f(x)$ сходится равномерно на $[a,b]$.

\end{lemma}

\begin{proofbox}{}

По критерию Коши
\[
\forall\varepsilon>0\ \exists n_0\
\forall n\geq n_0\
\forall p\in\mathbb{N}\
\forall x\in[a,b]:
\quad
|f_{n+p}(x)-f_n(x)|<\varepsilon
\]

Пусть $x\to a+0$. Так как функции $f_n$ непрерывны,
\[
f_{n+p}(x)\to f_{n+p}(a),
\qquad
f_n(x)\to f_n(a)
\]
Следовательно,
\[
|f_{n+p}(a)-f_n(a)|\leq\varepsilon
\]

Таким образом, предельная функция непрерывна на $[a,b]$.

\end{proofbox}

\newpage

\section{Свойства равномерной сходимости}

Пусть
\[
f_1(x),\, f_2(x),\,\ldots,\,f_n(x),\,\ldots
\overset{D}{\rightrightarrows} f(x)
\]
и при $x\to x_0$
\[
f_1(x)\to a_1,\qquad
f_2(x)\to a_2,\qquad
\ldots,\qquad
f_n(x)\to a_n,\qquad
\ldots
\]
а $f(x)\to A$.

\begin{theorem}{о перестановке пределов}
Пусть
\[
f_n:X\to\mathbb{R},\qquad D\subset X,\qquad
f_n(x)\overset{D}{\rightrightarrows}f(x)
\]
$x_0$ --- предельная точка множества $D$, и для любого $n\in\mathbb{N}$
\[
\lim_{x\to x_0}f_n(x)=a_n\in\mathbb{R}
\]

Тогда:

1) существует $\displaystyle\lim_{n\to\infty}a_n\in\mathbb{R}$;

2) $\lim_{x\to x_0}f(x) = \lim_{n\to\infty}a_n$
\end{theorem}

\begin{proofbox}{}
1) По критерию Коши: для любого $\varepsilon>0$ существует $n_0$ такое, что для любых $n\geq n_0$, $p\in\mathbb{N}$ и $x\in D$
\[
|f_{n+p}(x)-f_n(x)|<\varepsilon
\]

Переходя к пределу при $x\to x_0$, получаем
\[
|a_{n+p}-a_n|\leq\varepsilon
\]
Следовательно, последовательность $\{a_n\}$ фундаментальна, а значит, сходится.

Обозначим
\[
A=\lim_{n\to\infty}a_n
\]

2) Докажем, что
\[
A=\lim_{x\to x_0}f(x)
\]

Для любого $\varepsilon>0 \ \exists \delta > 0: x \in \overset{0}{U}_{\delta(x_0)} \Longrightarrow |f(x) - A| < \varepsilon$

Рассмотрим

\[
|f(x) - A| \leq |f(x) - f_n(x)| + |f_n(x) - a_n| + |a_n - A|<\varepsilon
\]

Следовательно,
\[
\lim_{x\to x_0}f(x)=A
\]
\end{proofbox}

\begin{corollary}
Если ряд
\[
\sum_{n=1}^{\infty}f_n(x)
\]
сходится равномерно на $D$ и существуют пределы $\displaystyle\lim_{x\to x_0}f_n(x)$, то
\[
\lim_{x\to x_0}\sum_{n=1}^{\infty}f_n(x)
=
\sum_{n=1}^{\infty}\lim_{x\to x_0}f_n(x)
\]
\end{corollary}

\begin{example}
\[
\lim_{x\to1-0}\lim_{n\to\infty}x^n
=
\lim_{x\to1-0}0
=
0
\]
но
\[
\lim_{n\to\infty}\lim_{x\to1-0}x^n
=
\lim_{n\to\infty}1
=
1
\]

Следовательно,
\[
x^n\not\rightrightarrows 0
\qquad\text{на }(0,1)
\]
\end{example}

\begin{theorem}{о непрерывности предельной функции}
Пусть
\[
f_n:X\to\mathbb{R},\qquad D\subset X,\qquad
f_n(x)\overset{D}{\rightrightarrows}f(x),
\qquad x_0\in D
\]
и $f_n$ непрерывна в точке $x_0$ для любого $n$.

Тогда $f$ непрерывна в точке $x_0$.
\end{theorem}

\begin{proofbox}{}
1) Если $x_0$ не является предельной точкой множества $D$, то $x_0$ изолирована, следовательно, $f$ непрерывна в $x_0$.

2) Пусть $x_0$ --- предельная точка множества $D$. Тогда по теореме о перестановке пределов
\[
\lim_{x\to x_0}f(x)
=
\lim_{n\to\infty}\lim_{x\to x_0}f_n(x)
=
\lim_{n\to\infty}f_n(x_0)
=
f(x_0)
\]
Следовательно, $f$ непрерывна в точке $x_0$.
\end{proofbox}

\begin{corollary}
Если $f_n$ непрерывна в точке $x_0$ для любого $n$ и ряд $\sum_{n=1}^{\infty}f_n(x)$ сходится равномерно, то его сумма
\[
S(x)=\sum_{n=1}^{\infty}f_n(x)
\]
непрерывна в точке $x_0$.
\end{corollary}

\begin{theorem}{об интегрировании предельной функции}
Пусть
\[
f_n:X\to\mathbb{R},\qquad [a,b]\subset X,\qquad
f_n\in\mathcal{R}[a,b]\quad\forall n
\]
и
\[
f_n\overset{[a,b]}{\rightrightarrows}f
\]

Тогда:

1) $f\in\mathcal{R}[a,b]$;

2)
\[
\int_a^x f_n(t)\,dt
\overset{[a,b]}{\rightrightarrows}
\int_a^x f(t)\,dt
\]
\end{theorem}

\begin{proofbox}{}
1) Критерий интегрируемости:
\[
f\in\mathcal{R}[a,b]
\Longleftrightarrow
\forall\varepsilon>0\ \exists\delta>0\ \forall\tau:
\lambda(\tau)<\delta
\quad\Longrightarrow\quad
S_\tau-s_\tau<\varepsilon
\]
где
\[
S_\tau-s_\tau=\sum \omega_i\Delta x_i
\]

\[
|f(x)-f(y)|
\leq
|f(x)-f_n(x)|
+
|f_n(x)-f_n(y)|
+
|f_n(y)-f(y)|
\]

Следовательно,
\[
\omega_i(f)\leq 2\varepsilon+\omega_i(f_n)
\]
и
\[
\Omega_f
\leq
2\varepsilon(b-a)+\Omega_{f_n}
<
\varepsilon\bigl(2(b-a)+1\bigr)
\quad\Longrightarrow\quad
f\in\mathcal{R}[a,b]
\]

2)
\[
\left|
\int_a^x f_n-\int_a^x f
\right|
\leq
\int_a^x|f_n-f|\,dt
<
\frac{\varepsilon}{b-a}(x-a)
\leq\varepsilon
\]
\end{proofbox}

\begin{corollary}
Если $f_n\in\mathcal{R}[a,b]$ и ряд
\[
\sum_{n=1}^{\infty}f_n(x)
\]
сходится равномерно на $[a,b]$, то
\[
\int_a^x\sum_{n=1}^{\infty}f_n(t)\,dt
=
\sum_{n=1}^{\infty}\int_a^x f_n(t)\,dt
\]
\end{corollary}

\begin{theorem}{о дифференцировании предельной функции}
Пусть
\[
f_n\in C^1[a,b],\qquad x_0\in[a,b]
\]
последовательность $\{f_n(x_0)\}$ сходится и
\[
f_n'(x)\overset{[a,b]}{\rightrightarrows}g(x)
\]

Тогда:

1)
\[
f_n(x)\overset{[a,b]}{\rightrightarrows}f(x)
\]

2)
\[
f\in C^1[a,b],
\qquad
f'(x)=g(x)
\]
\end{theorem}

\begin{proofbox}{}
1) Так как
\[
f_n' \overset{[x_0,x]}{\rightrightarrows} g
\]
то
\[
\int_{x_0}^{x} f_n' \overset{[x_0,x]}{\rightrightarrows} \int_{x_0}^{x} g
\]

\[
\lim_{n\to\infty}\bigl(f_n(x)-f_n(x_0)\bigr)
=
\int_{x_0}^{x} g\,dt
\quad\Longrightarrow\quad
f_n(x)\longrightarrow C+\int_{x_0}^{x}g(t)\,dt=f(x)
\]
где
\[
\lim_{n\to\infty}f_n(x_0)=C
\]

2) $f$ дифференцируема, и
\[
f'(x)
=
\left(C+\int_{x_0}^{x}g(t)\,dt\right)'
=
g(x)
\]
(по теореме Барроу), $g(x)$ непрерывна, так как $f_n'$ непрерывны.

Следовательно,
\[
f\in C^1
\]
\end{proofbox}

\begin{example}

\textbf{Применение}

Найти
\[
S(x)=\sum_{n=1}^{\infty}nx^n=x+2x^2+3x^3+\ldots
\]

Ряд сходится на $(-1,1)$ неравномерно и сходится равномерно на $[-a,a]\subset(-1,1)$.

\[
(x^n)'=nx^{n-1}
\]

I способ:
\[
S(x)
=
x\sum_{n=1}^{\infty}nx^{n-1}
=
x\sum_{n=1}^{\infty}(x^n)'
=
x\left(\sum_{n=1}^{\infty}x^n\right)'
=
x\left(\frac{x}{1-x}\right)'
=
\frac{x}{(1-x)^2}
\]

II способ:

$S(x)$ сходится равномерно на $[-a,a]$.

\[
S(x)
=
\sum_{n=1}^{\infty}nx^n
=
x\sum_{n=1}^{\infty}nx^{n-1}
=
S_1(x)
\]

\[
\int_0^x S_1(t)\,dt
=
\int_0^x\left(\sum_{n=1}^{\infty}nt^{n-1}\right)dt
=
\sum_{n=1}^{\infty}\int_0^x nt^{n-1}\,dt
=
\sum_{n=1}^{\infty}x^n
=
\frac{x}{1-x}
\]

Следовательно,
\[
S_1(x)
=
\left(\frac{x}{1-x}\right)'
\]

\end{example}

\newpage

\section{Степенные ряды}

\begin{definition}
$\sum_{n=0}^\infty a_n(x-x_0)^n = a_0 + a_1(x-x_0) + a_2(x-x_0)^2 + ...$ - \textbf{степенной ряд}

При $x = x_0$ - сходится

Замена $t = x - x_0: \qquad \sum_{n=0}^\infty a_nt^n$
\end{definition}

\begin{theorem}{ - Первая теорема Абеля}

1) Если ряд $\sum^\infty$ сходится при $x=x_1$, то он сходится $\forall x: |x|<|x_1|$

2)  Если при $x=x_2$ ряд расходится, то он расходится $\forall x: |x|>|x_2|$

\end{theorem}

\begin{proofbox}{}
Пусть ряд $\sum_{n=0}^\infty a_n x_1^n$ сходится и $|x|<|x_1|$

Оценим: 
\[
|a_n x_n| = |a_n x_1^n| \cdot \left| \frac{x^n}{x_1^n} \right|
\]

Так как $|a_n x_1^n| \longrightarrow 0$, то оно ограничено. Тогда 

\[
|a_n x_1^n| \cdot \left| \frac{x^n}{x_1^n} \right|
\leq c \ \cdot \left| \frac{x^n}{x_1^n} \right|
\]

Но 

\[
\left| \frac{x}{x_1} \right| < 1
\]

Следовательно, ряд сходится абсолютно.

2) Пусть при $x: |x|>|x_2|$ ряд сходится. Тогда по п.1 в точке $x_2$ ряд сходится. Противоречие
\end{proofbox}

\begin{corollary}
$\exists! R \in [0, +\infty]: $

на $(-R, R)$ ряд сходится,

а на $(-\infty, -R)\cup(R, +\infty)$ ряд расходится

В таком случае:

R - радиус сходимости

$(-R, R)$ - интервал сходимости

\end{corollary}

\begin{example}
\[
\text{1) } \sum_{n=0}^\infty x^n \ D = (-1, 1), \ R = 1
\]

\[
\text{2) }\sum_{n=0}^\infty \frac{x_n}{n^2} \ \ D = [-1, 1], \ R = 1
\]

\[
\text{3) }\sum_{n=0}^\infty \frac{x_n}{n} \ \ D = [-1, 1), \ R = 1
\]

\[
\text{4) }\sum_{n=0}^\infty n! \cdot x^n \ \ D = \{0\}, \ R = 0
\]

\[
\text{5) }\sum_{n=0}^\infty \frac{x_n}{n!} \ \ D = \mathbb{R}, \ R = +\infty
\]

\end{example}

\begin{theorem}{- Формула Коши-Адамара}

Для ряда $\sum_{n=0}^\infty a_nx^n$ верно, что:

\[
R = \frac{1}{\overline{\lim_{n\to\infty}}\sqrt[n]{|a_n|}}
\]

\end{theorem}

\begin{proofbox}{}

Применим радикальный признак Коши:

\[
l = \overline{\lim_{n\to\infty}}\sqrt[n]{|a_n x^n|}
\]

Если $l < 1 \Longrightarrow$ ряд $\sum|a_nx^n|$ сходится

Если $l > 1 \Longrightarrow$ ряд $\sum|a_nx^n|$ расходится, так как не стремится к 0, а значит $\sum a_nx^n$ расходится.

\[
l = |x|\overline{\lim_{n\to\infty}}\sqrt[n]{|a_n|} < 1 \Longleftrightarrow |x| < \frac{1}{\overline{\lim_{n\to\infty}}\sqrt[n]{|a_n|}}
\]

\end{proofbox}

\begin{note}
Форула Коши-Адамара верна в $\overline{\mathbb{R}}$, то есть

\[
\overline{\lim} = 0 \Longrightarrow R = + \infty
\]
и 
\[
\overline{\lim} = +\infty \Longrightarrow R = 0
\]

\end{note}

\begin{theorem}{о равномерной сходимости степенного ряда}

Ряд $\sum_{n=0}^\infty a_nx^n$ сходится равномерно $\forall [a, b] \subset (-R, R)$ (на любом отрезке внутри интервала сходимости)

\end{theorem}

\begin{proofbox}{}
\[
|a_n x^n| \leq |a_n| \cdot c^n
\]

Где $c = max\{|a|, |b|\}$

А такой ряд сходится. Следовательно исходный ряд сходится равномерно по признаку Вейерштрасса

\end{proofbox}

\begin{corollary}

1) Сумма степенного ряда непрерывна на $(-R, R)$

2) На $[a, b] \subset (-R, R)$ можно интегрировать

3) 

\end{corollary}

\begin{lemma}{о радиусах сходимости}

Ряды
\[
\sum^\infty a_n x^n, \qquad \sum^\infty a_n n x^{n-1}, \qquad \sum^\infty \frac{a_n}{n+1} x^{n+1}
\]

Имеют равный радиус сходимости

\end{lemma}

\begin{proofbox}{}

По формуле Коши-Адамара для $\sum^\infty a_n x^n$

\[
R_1 = \frac{1}{\overline{\lim\sqrt[n]{|a_n|}}}
\]

Для $\sum^\infty a_n n x^{n-1} \text{ и } \sum^\infty \frac{a_n}{n+1} x^{n+1}:$

\[
R_2 = \frac{1}{\overline{\lim\sqrt[n]{|a_n \cdot n|}}} \longrightarrow R_1
\\
R_3 = \frac{1}{\overline{\lim\sqrt[n]{|\frac{a_n}{n+1}|}}} \longrightarrow R_1
\]

Так как $\sqrt[n]{n} \text{ и } \sqrt[n]{n+1} \to 1$

\end{proofbox}

\begin{corollary}

\[
S(x) =
\left(
\sum_{n=0}^\infty a_n x^n
\right)' =
\sum_{n=1}^\infty a_n n x^{n-1}
\\
S^{(m)}(x) = \sum_{n=m}^\infty a_n \frac{n!}{(n-m)!} x^{n-m}
\]

\end{corollary}

\begin{theorem}{- Вторая теорема Абеля}

Пусть ряд $\sum_{n=0}^\infty a_n x^n$ имеет радиус сходимости $R \in (0, +\infty)$ и при $x = R \ (x = -R)$ ряд сходится. Тогда он сходится равномерно на $[0, R] ([-R, 0])$

\end{theorem}

\begin{proofbox}{}

Представим

\[
a_n x^n = a_n R^n \cdot \left(\frac{x}{R}\right)^n
\\
\sum^\infty a_nR^n \ \text{сходятся равномерно}
\\
\left(\frac{x}{R}\right)^n \downarrow \forall x \in [0, R] \text{ и } \left|\left(\frac{x}{R}\right)^n\right| \leq 1
\]

Следовательно по признаку Абеля-Дирихле ряд сходится равномерно

\end{proofbox}

\begin{theorem}{о единственности степенного ряда}

Пусть $\forall x \in (x_0 - R, x_0 + R)$:

\[
f(x) = \sum_{n=0}^\infty a_n (x - x_0)^n
\]

Тогда

\[
a_n = \frac{f^{(n)}(x_0)}{n!}
\]

То есть это ряд Тейлора

\end{theorem}

\begin{proofbox}{}

\[
f^{(m)}(x) = \sum_{n=m}^\infty a_n \frac{n!}{(n-m)!} (x - x_0)^{n-m}
\]

При $x = x_0:$

\[
f^{(m)}(x) = a_m \cdot m!
\]

\end{proofbox}

\newpage

\section{Ряд Тейлора}

\subsection{Определение}

\begin{definition}

Пусть $f \in C^\infty(x_0)$. Тогда ряд

\[
\sum_{n=0}^\infty a_n(x - x_0)^n \text{, где } 
a_n = \frac{f^{(n)}(x_0)}{n!}
\]

Называется \textbf{рядом Тейлора} функции f с центром $x_0$

\end{definition}

\begin{example}

\begin{flalign*}
f(x) &=
\begin{cases}
e^\frac{-1}{x^2}, & x \not=0,\\
0, & x=0
\end{cases}
&
\end{flalign*}

\[
f'(0) = \lim_{\Delta x \to 0}\frac{f(0+\Delta x) - f(0)}{\Delta x} = \lim_{\Delta x \to 0}\frac{1}{\Delta x} e^{\frac{-1}{x^2}} = 0
\]
$x\not= 0:$
\[
f'(x) = (e^\frac{-1}{x^2})' = \frac{2}{x^3}e^\frac{-1}{x^2}
\]

\[
f''(0) = \lim_{\Delta x \to 0}\frac{2}{\Delta x^4} e^{\frac{-1}{x^2}} = 0
\]
$x\not= 0:$
\[
f'(x) = (e^\frac{-1}{x^2})'' = (\frac{2}{x^3}e^\frac{-1}{x^2})' = (-\frac{6}{x^4} + \frac{4}{x^6})\cdot e^\frac{-1}{x^2}
\]

И так далее. Получается:

\[
f^{(m)}(0) = 0 \qquad f(x) \sim \sum^\infty_{n=0} 0
\]

Получается, что функцию f нельзя приблизить степенной функцией

\end{example}

\begin{definition}

Пусть

\[
f(x) \sim \sum^\infty_{n=0}a_n(x - x_0)^n \ - \text{ряд Тейлора}
\]

Тогда пусть $P_n(x) = \sum^\infty_{k=0}a_k(x - x_0)^k \ - $ \textbf{частичная сумма}

$R_n = f(x) - P_n(x) \ -$ \textbf{остаток}

$f(x) = P_n(x) + R_n(x)$

\end{definition}

\begin{lemma}

\[
\sum^\infty_{n=0}a_n(x - x_0)^n = f(x) \Longleftrightarrow
R_n(x) \longrightarrow 0
\]

\end{lemma}

\begin{theorem}{о достаточном условии сходимости ряда Тейлора к f(x)}

Если на $[x, x_0]$ или $[x_0, x] \ |f^{(m)}(x)| \leq c$, то ряд Тейлора сходится к f(x)

\end{theorem}

\begin{proofbox}{}

$f(x) = P_n(x) + R_n(x)$, где

\[
R_n(x) = \frac{f^{(n+1)}(\xi)}{(n+1)!}(x - x_0)^{(n+1)}, \qquad \xi\in [x, x_0]
\\ \Longrightarrow \left|R_n(x)\right| \leq \frac{c |x - x_0|^{n+1}}{(n+1)!} \longrightarrow 0
\]

Так как

\[
\frac{|x - x_0|^{n+1}}{(n+1)!} \longrightarrow 0
\]

\end{proofbox}

\subsection{Ряды Маклорена элементарных функций}

\begin{example}

1) $f(x) = e^x$

\[
f^{(n)}(x) = e^x (= 1 \text{ при } x = 0)
\\
e^x \sim 1 + \frac{x}{1!} + \frac{x^2}{2!} + \frac{x^3}{3!} + ... = \sum^\infty_{n=0} \frac{x^n}{n!}
\]

Пусть $x \in [-r, r]$. Тогда $|f^{(n)}(x)| = e^x \leq e^r = c$. Следовательно

\[
e^x = \sum_{n=0}^\infty \frac{x^n}{n!} \ \forall x \in \mathbb{R}
\]

2) $f(x) = \sin(x)$

\[
\sin x = x - \frac{x^3}{3!} + \frac{x^5}{5!} - ... = \sum^\infty_{n=0} \frac{(-1)^n \cdot x^{2n+1}}{(2n+1)!}
\]
Так как $|(\sin x)^{(n)}| \leq 1$

3) $f(x) = \cos(x)$

\[
\cos x = 1 - \frac{x^2}{2!} + \frac{x^4}{4!} - ... = \sum^\infty_{n=0} \frac{(-1)^n \cdot x^{2n}}{(2n)!}
\]

Аналогично синусу

4) $f(x) = (1+x)^\alpha$

\[
(1+x)^\alpha \sim 1 + \alpha x + \frac{\alpha(\alpha - 1)}{2!} \cdot x^2 + ... = \sum_{n=0}^\infty \frac{\alpha!}{(\alpha - n + 2)! \cdot n!} \cdot x^n
\]

Докажем равенство. Пусть

\[
S(x) = \sum_{n=0}^\infty \frac{\alpha!}{(\alpha - n + 2)! \cdot n!} \cdot x^n
\]

\[
S'(x) = \sum_{n=1}^\infty \frac{\alpha!}{(\alpha - n + 2)! \cdot (n-1)!} \cdot x^{n-1}
\\
= \sum_{n=0}^\infty \frac{\alpha!}{(\alpha - n + 1)! \cdot n!} \cdot x^n
\]

Рассмотрим

\[
xS' + S'
=
\sum_{k=1}^{\infty}
\frac{\alpha\cdots(\alpha-k+1)\,k}{k!}x^k
+
\sum_{k=0}^{\infty}
\frac{\alpha\cdots(\alpha-k+1)(\alpha-k)}{k!}x^k
\]

\[
=
\alpha
+
\underbrace{
\sum_{k=1}^{\infty}
\frac{\alpha\cdots(\alpha-k+1)}{k!}x^k
}_{S-1}
\bigl(k+\alpha-k\bigr)
=
\alpha+\alpha(S-1)
\]
Получается
\[
(x+1)S'=\alpha S
\]

То есть

\[
(1+x)^\alpha \sim
\underbrace{
\sum_{k=0}^{\infty}
\frac{\alpha(\alpha-1)(\alpha-2)\cdots(\alpha-k+1)}{k!}\,x^k
}_{S(x)}
\]

\[
x\in(-1,1),\qquad \forall\alpha.
\]

\[
S'\cdot x+S'=\alpha\cdot S,
\qquad
S'=\frac{dS}{dx}.
\]

\[
(1+x)\cdot dS=\alpha\cdot S\cdot dx.
\]

\[
\int\frac{dS}{S}
=
\int\frac{\alpha\cdot dx}{1+x}.
\]

\[
d\ln|S|=d\bigl(\alpha\ln|1+x|\bigr).
\]

\[
\ln|S|=\alpha\cdot\ln|1+x|+C.
\]

Знаем $S(0)=1\Rightarrow C=0$.

\[
\Rightarrow\quad \ln|S|=\alpha\cdot\ln|1+x|.
\]

\[
|S|=|1+x|^\alpha,
\qquad 1+x>0.
\]

\[
S=(1+x)^\alpha.
\]

Например: 
\[
\frac{1}{1-x}
=
\sum_{k=0}^{\infty}
\frac{(-1)^k\cdot k!\cdot(-x)^k}{k!}
=
\sum_{k=0}^{\infty}x^k,
\qquad x\in(-1,1).
\]

\[
\frac{1}{1+x}
=
\sum_{k=0}^{\infty}(-1)^k x^k,
\qquad x\in(-1,1).
\]

\

5) $f(x) = \ln(1 + x)$

Так как

\[
\frac{1}{1 + x} = \sum_{k=0}^\infty (-1)^k x^k \text{ сходится равномерно на } [0, x], то
\\
\ln(1+x) = \int_0^x \frac{dt}{1+t} = \sum_{k=0}^\infty (-1)^k t^k dt =
\\
\sum_{k=0}^\infty \frac{(-1)^k \cdot t^{k+1}}{k+1} = \sum_{k=1}^\infty \frac{(-1)^{k-1} \cdot t^k}{k}
\]

Сходится на $(-1, 1] \Longrightarrow$ по 2-ой теореме Абеля сходится равномерно на $[0, 1]$

\end{example}

\newpage

\section{Сходимость почти всюду и по мере}

Пусть $(X, \mathcal{A}, \mu)$ - измеримое пространство

\begin{definition}

Утверждение P(x) истинно на $A \in \mathcal{A}$ почти всюду, если
\[
\exists C \in \mathcal{A}, C \subset A \text{ и}
\\
P(x) \text{ истинно } \forall x \in C \text{ и } \mu(A\backslash C) = 0
\]

\end{definition}

\begin{example}
1) \textbf{Функция Дирихле} - $D(x) = 0$ почти всюду по мере Лебега на $\mathbb{R}$

2) $x^n \longrightarrow 0$ почти всюду по мере Лебега на $[-1, 1]$

\end{example}

\begin{note}

\textbf{О единственности предела почти наверное}

Если 

\[
f_n \xrightarrow{п.в. на E} f(x) \text{ и}
\\
f_n \xrightarrow{п.в. на E} g(x), \text{ то}
\]

$f = g$ почти всюду на E $(f \sim g)$

\end{note}

\begin{definition}

Пусть $f_n, f -$ измеримы и почти всюду конечны на E

\[
f_n \xrightarrow{\mu} f \text{ на } E
\\
\Longleftrightarrow \forall \varepsilon > 0 \qquad \mu E(|f_n - f| \geq \varepsilon) \longrightarrow 0
\]

\end{definition}

\begin{example}

1) $x^n \xrightarrow{\lambda} 0$ почти всюду по мере Лебега на $[-1, 1]$

\[
|x^n| \geq \varepsilon \qquad |x| \geq \sqrt[n]\varepsilon\,\ n \to \infty \Longrightarrow \mu E \to 0
\]

2) $f_n(x) = \chi_{[n , n + 1]}(x) \to 0$ на $\mathbb{R}$

\[
|f_n(x)| \geq \varepsilon \qquad x \in [n, n + 1]
\\
\lambda[n, n + 1] = 1 \not\to 0 \Longrightarrow f_n \not \xrightarrow{\lambda} 0
\]

3) \textbf{Пример Рисса}

Пусть 
\[
\begin{aligned}
f_0&=\chi_{[0,1]},\\
f_1&=\chi_{\left[0,\frac12\right]},\\
f_2&=\chi_{\left[\frac12,1\right]},\\
f_3&=\chi_{\left[0,\frac14\right]},\\
&\vdots
\end{aligned}
\]

Тогда

\[
f_k\xrightarrow{\lambda}0.
\]
\[
f_k\not\to0
\qquad \forall x\in[0,1].
\]

\end{example}

\end{document}